% =====================================================================
%  รายงานการทดสอบ — ตารางสถานการณ์ทดสอบ / Test Case
%  Compile: XeLaTeX  (Overleaf: Menu > Settings > Compiler > XeLaTeX)
%
%  โฟลเดอร์นี้แยกจากเล่มคู่มือ ใช้เก็บ test case ตามรูปแบบตารางในเล่มรายงาน
%  วิธีใช้: กรอกข้อมูลส่วนที่ 1 แล้วเพิ่มไฟล์ใน scenarios/ แล้ว \input ในส่วนที่ 3
% =====================================================================
\documentclass[11pt,a4paper]{article}
\usepackage{testreport}

% ---------------------------------------------------------------------
% 1) ข้อมูลเล่ม
% ---------------------------------------------------------------------
\projectCode{CS 2568/01}
\titleTH{ระบบดูแลและคัดกรองสุขภาพจิตนิสิต (NUCARE)}
\titleEN{NUCARE — Student Mental Health Care and Screening System}
\doctype{รายงานการทดสอบ}

\systemUnderTest{ระบบเว็บแอปพลิเคชัน ``NUCARE''}
\systemVersion{1.0}
\testerName{นายสมชาย ใจดี}
\testDateTH{วันที่ 22 ตุลาคม พ.ศ.\ 2568}

\firstStudentID{653450001-1}
\firstStudentPrefixTH{นาย}
\firstStudentNameTH{สมชาย ใจดี}
\secondStudentID{653450002-9}
\secondStudentPrefixTH{นางสาว}
\secondStudentNameTH{สมหญิง รักเรียน}

\advisorTH{ผศ.ดร.ประสิทธิ์ สร้างสรรค์}
\courseCode{CP354 772}
\courseName{โครงงานคอมพิวเตอร์ 2}
\semester{2}
\academicYearTH{2568}
\submitMonthTH{ตุลาคม}
\submitYearTH{2568}

\begin{document}

% ---------------------------------------------------------------------
% 2) ปก (ลบ \maketestcover ถ้าไม่ต้องการปก)
% ---------------------------------------------------------------------
\maketestcover

\section*{สรุปการทดสอบ}
เอกสารฉบับนี้รวบรวมสถานการณ์ทดสอบ (Test Scenario) และกรณีทดสอบ (Test Case)
ของระบบเว็บแอปพลิเคชัน ``NUCARE'' เพื่อบันทึกขั้นตอน ผลลัพธ์ที่คาดหวัง ผลลัพธ์จริง และผลการทดสอบ
ในรูปแบบตารางมาตรฐานของเล่มรายงานโครงงาน

\begin{center}
\begin{tabularx}{0.92\linewidth}{|c|Y|c|c|}
\hline
\rowcolor{tcHeader}
\textbf{รหัส} & \textbf{ชื่อสถานการณ์ทดสอบ} & \textbf{จำนวนกรณี} & \textbf{ผลรวม} \\
\hline
TC01 & การเข้าสู่ระบบของผู้ใช้งานผ่านหน้าเว็บแอปพลิเคชัน & 4 & Pass \\
\hline
TC02 & การทำแบบประเมินสุขภาพจิตของนิสิตผ่านหน้าเว็บแอปพลิเคชัน & 3 & Pass \\
\hline
TC03 & การตรวจสอบสิทธิ์บทบาทเพิ่มเติมและการเข้าถึงข้อมูลนิสิต & 3 & Pass \\
\hline
TC04 & การบันทึกนัดหมายของทีมสุขภาพจิตและแสดงให้นิสิต & 2 & Pass \\
\hline
TC05 & การดูสรุปรายงานและทางลัดตั้งค่าของผู้ดูแลระบบ & 2 & Pass \\
\hline
TC06 & การอัปเดตสถานะ referral ของงานพัฒนานิสิต & 3 & Pass \\
\hline
TC07 & การรับและติดตามข้อสังเกตจากผู้สอนโดยอาจารย์ที่ปรึกษา & 2 & Pass \\
\hline
TC08 & การคุยกับเพื่อน AI และการกันถามการบ้าน & 3 & Pass \\
\hline
TC09 & การตรวจจับ Wellness Trend Indicators จากแชทนิสิต & 4 & Pass \\
\hline
TC10 & การตั้งค่าคลื่นแบบประเมินและกฎคัดเลือกกลุ่มเป้าหมาย & 3 & Pass \\
\hline
TC11 & การย่อชื่อ/อวตาร์โดยไม่ใช้ท้ายรหัสนิสิต & 3 & Pass \\
\hline
TC12 & การสรุปสถานะติดตามข้อสังเกตผู้สอนจากสถานะ Alert & 3 & Pass \\
\hline
TC13 & การยุบบทบาทผู้สอนและการแจ้งข้อสังเกตเฉพาะนอกสังกัด & 4 & Pass \\
\hline
TC14 & การสร้างเส้นทางติดตามข้อสังเกตจากเหตุการณ์จริง + ส่งต่อต่อได้ & 5 & Pass \\
\hline
TC15 & การตั้งความถี่แบบสอบถามแบบกำหนดเองและร่างคลื่นรอบถัดไป & 3 & Pass \\
\hline
TC16 & โครงส่งอีเมลแบบเปลี่ยน provider ได้และช่องทางแจ้งเตือน & 3 & Pass \\
\hline
TC17 & สัญญาณจากแชท AI ในประวัตินิสิต (ไม่ขึ้นกล่องแจ้งเตือน) & 3 & Pass \\
\hline
TC18 & นโยบายปิดแจ้งเตือน: dismiss / บังคับรับเรื่องก่อนปิดเคส & 3 & Pass \\
\hline
TC19 & แชทเพื่อน AI ผ่าน KKU Gen AI (Gemini/GPT/Claude) & 3 & Pass \\
\hline
TC20 & คำขอปรึกษานิสิต--ที่ปรึกษา และการนัดปรึกษาสุขภาพจิต & 6 & Pass \\
\hline
TC21 & เส้นทางหมุดติดตามการส่งต่อ และการเจรจานัดปรึกษาสองฝ่าย & 5 & Pass \\
\hline
TC22 & แยกคิวเจรจานัดของที่ปรึกษา และการเรียงนัดใกล้ถึง & 4 & Pass \\
\hline
TC23 & กล่องงานนิสิต: ต้องตอบ ข้อเสนอนัด และนัดใกล้ถึง & 6 & Pass \\
\hline
TC24 & กันนิสิตมีคำขอปรึกษาเปิดพร้อมกันหลายรอบ & 3 & Pass \\
\hline
TC25 & รูปแบบออนไลน์/พบตัว ยกเลิกโดยนิสิตหรืออาจารย์ (ป้ายเหลือง) และสถานที่จากอาจารย์ & 5 & Pass \\
\hline
TC26 & การแจ้งช่วยเหลือด่วน (SOS) เมื่อส่งไม่สำเร็จ & 3 & Pass \\
\hline
TC27 & การแบ่งหน้ารายการยาวแบบเดียวกับหน้า admin & 3 & Pass \\
\hline
TC28 & ปรับพอร์ทัลผู้บริหารให้ตรง PDPA และการใช้งาน & 5 & Pass \\
\hline
TC29 & Dossier นิสิตผู้บริหารและความกว้างเลย์เอาต์ & 4 & Pass \\
\hline
TC30 & เวิร์กโฟลว์ผู้เชี่ยวชาญสุขภาพจิต (รับคิว นัด บันทึกผล ประวัติ) & 6 & Pass \\
\hline
TC31 & เสริมความพร้อมใช้งานจริงของพอร์ทัลผู้เชี่ยวชาญสุขภาพจิต & 5 & Pass \\
\hline
TC32 & การ hardening พอร์ทัลงานพัฒนานิสิตให้ใช้งานจริง & 4 & Pass \\
\hline
TC33 & การแจ้งช่วยเหลือด่วน (SOS) มีผู้รับและดำเนินการได้ & 4 & Pass \\
\hline
TC34 & ผลคัดกรองความเสี่ยงสูงแจ้งที่ปรึกษาทางแอปและอีเมล & 3 & Pass \\
\hline
TC35 & การเชื่อมโยงบทบาทหลังแก้จุดขาด handoff & 4 & Pass \\
\hline
TC36 & การจัดการภาพประชาสัมพันธ์โดยแอดมิน & 3 & Pass \\
\hline
TC37 & Toast หลังกดปุ่ม แยกข้อความตามการกระทำ + helper คัด alerts & 5 & Pass \\
\hline
TC38 & สิทธิ์งานพัฒนานิสิตจัดการภาพประชาสัมพันธ์หน้านิสิต & 3 & Pass \\
\hline
TC39 & การ hardening ความพร้อมใช้งานจริง (รหัสผ่าน อีเมล แชท) & 4 & Pass \\
\hline
TC40 & นิสิตขอนัดผู้เชี่ยวชาญสุขภาพจิตผ่านช่องทางเฉพาะ & 3 & Pass \\
\hline
TC41 & แชทเพื่อน AI จำต่อเนื่องแบบประหยัดและกันทำการบ้าน & 5 & Pass \\
\hline
TC42 & แยกคิวรอรับกับเคสที่ดูแล และปรับ UX triage & 6 & Pass \\
\hline
TC43 & เปิดโมดัลแล้วแถบบนไม่หาย / ไม่ดันเลย์เอาต์ด้วยสกอลล์บาร์ซ้อน & 3 & Pass \\
\hline
TC44 & ปฏิทินนัดใกล้ถึง 7 วันบนหน้าแรกที่ปรึกษา & 3 & Pass \\
\hline
TC45 & จำกัดรายการแจ้งเตือน ปฏิทินหลายวัน และกันชนนัด 1 ชั่วโมง & 4 & Pass \\
\hline
TC46 & การตั้งค่า Secure flag ของ cookie ยืนยันตัวตน & 4 & Pass \\
\hline
\end{tabularx}
\end{center}

\clearpage

% ---------------------------------------------------------------------
% 3) ตารางสถานการณ์ทดสอบ — เพิ่ม \input{scenarios/...} ตามจำนวนเคส
% ---------------------------------------------------------------------
\section*{รายละเอียดสถานการณ์ทดสอบ}
% TC01 — การเข้าสู่ระบบ (รูปแบบตารางตามเล่มอ้างอิง รวมหน้าต่อ)
\begin{testscenario}{1}{1}
  \tcmeta{TC01}{CS 2568/01}{ระบบเว็บแอปพลิเคชัน ``NUCARE''}{นายสมชาย ใจดี}{วันที่ 22 ตุลาคม พ.ศ.\ 2568}{1.0}
  \tcname{การเข้าสู่ระบบของผู้ใช้งานผ่านหน้าเว็บแอปพลิเคชัน}
  \tcpre{ผู้ใช้งานมีบัญชีในระบบแล้ว และสามารถเปิดหน้าเข้าสู่ระบบผ่านเว็บเบราว์เซอร์ได้}
  \tcdesc{ทดสอบการทำงานของระบบเข้าสู่ระบบ โดยกรอกข้อมูลในแต่ละกรณีให้ครอบคลุมทั้งข้อมูลถูกต้องและข้อมูลไม่ถูกต้อง
    เพื่อยืนยันว่าระบบสามารถตรวจสอบสิทธิ์ แสดงข้อความแจ้งเตือน หรือนำผู้ใช้ไปยังแดชบอร์ดได้ถูกต้องตามเงื่อนไขที่ออกแบบไว้}

  \tccase{1}{ตรวจสอบการเข้าสู่ระบบเมื่อผู้ใช้กรอกข้อมูลถูกต้องครบถ้วน}{%
    \begin{tcsteps}
      \item เปิดหน้า ``เข้าสู่ระบบ'' ของระบบ
      \item กรอกชื่อผู้ใช้ = student01
      \item กรอกรหัสผ่าน = Password123!
      \item กดปุ่ม ``เข้าสู่ระบบ''
    \end{tcsteps}}{%
    ระบบตรวจสอบสิทธิ์สำเร็จ และนำผู้ใช้ไปยังแดชบอร์ดตามบทบาท}{%
    ระบบตรวจสอบสิทธิ์สำเร็จ และนำผู้ใช้ไปยังแดชบอร์ดตามบทบาท}{Pass}

  \tccase{2}{ตรวจสอบเมื่อผู้ใช้กรอกชื่อผู้ใช้ว่าง หรือมีความยาวน้อยกว่า 4 ตัวอักษร}{%
    \begin{tcsteps}
      \item เปิดหน้า ``เข้าสู่ระบบ'' ของระบบ
      \item กรอกชื่อผู้ใช้ = ว่าง หรือ stu
      \item กรอกรหัสผ่าน = Password123!
      \item กดปุ่ม ``เข้าสู่ระบบ''
    \end{tcsteps}}{%
    ใต้ช่องชื่อผู้ใช้แสดงข้อความ ``ชื่อผู้ใช้ต้องมีอย่างน้อย 4 ตัวอักษร''}{%
    ใต้ช่องชื่อผู้ใช้แสดงข้อความ ``ชื่อผู้ใช้ต้องมีอย่างน้อย 4 ตัวอักษร''}{Pass}

  \tccasebreak
  \tccasecont{2 (ต่อ)}{%
    \begin{tcsteps}
      \setcounter{enumi}{4}
      \item กรอกยืนยันว่ารหัสผ่านยังเป็น Password123!
      \item กดปุ่ม ``เข้าสู่ระบบ'' อีกครั้งหลังแก้ชื่อผู้ใช้สั้น
    \end{tcsteps}}

  \tccase{3}{ตรวจสอบเมื่อผู้ใช้กรอกรหัสผ่านว่าง หรือรหัสผ่านไม่ถูกต้อง}{%
    \begin{tcsteps}
      \item เปิดหน้า ``เข้าสู่ระบบ'' ของระบบ
      \item กรอกชื่อผู้ใช้ = student01
      \item กรอกรหัสผ่าน = ว่าง หรือ wrong-password
      \item กดปุ่ม ``เข้าสู่ระบบ''
    \end{tcsteps}}{%
    ระบบปฏิเสธการเข้าสู่ระบบ และแสดงข้อความแจ้งว่าชื่อผู้ใช้หรือรหัสผ่านไม่ถูกต้อง}{%
    ระบบปฏิเสธการเข้าสู่ระบบ และแสดงข้อความแจ้งว่าชื่อผู้ใช้หรือรหัสผ่านไม่ถูกต้อง}{Pass}

  \tccase{4}{ตรวจสอบเมื่อผู้ใช้กรอกรหัสผ่านมีความยาวน้อยกว่า 8 ตัวอักษร}{%
    \begin{tcsteps}
      \item เปิดหน้า ``เข้าสู่ระบบ'' ของระบบ
      \item กรอกชื่อผู้ใช้ = student01
      \item กรอกรหัสผ่าน = Pass12
      \item กดปุ่ม ``เข้าสู่ระบบ''
    \end{tcsteps}}{%
    ใต้ช่องรหัสผ่านแสดงข้อความ ``รหัสผ่านต้องมีอย่างน้อย 8 ตัวอักษร''}{%
    ใต้ช่องรหัสผ่านแสดงข้อความ ``รหัสผ่านต้องมีอย่างน้อย 8 ตัวอักษร''}{Pass}
\end{testscenario}

% TC02 — แบบประเมินสุขภาพจิตของนิสิต (แสดงเกณฑ์ก่อน/หลัง และผลหลังส่ง)
\begin{testscenario}{2}{3}
  \tcmeta{TC02}{CS 2568/01}{ระบบเว็บแอปพลิเคชัน ``NUCARE''}{นายสมชาย ใจดี}{วันที่ 5 พฤศจิกายน พ.ศ.\ 2569}{1.2}
  \tcname{การทำแบบประเมินสุขภาพจิตของนิสิตผ่านหน้าเว็บแอปพลิเคชัน}
  \tcpre{ผู้ใช้งานเข้าสู่ระบบด้วยบทบาทนิสิตแล้ว และมีคลื่นแบบประเมินที่เปิดอยู่พร้อม assignment}
  \tcdesc{ทดสอบการส่งแบบประเมินเมื่อนิสิตอบครบทุกข้อ ยืนยันว่าบันทึกคำตอบได้
    และนิสิตเห็นเกณฑ์ก่อนทำ รวมถึงคะแนน/แปลผลหลังส่ง โดย options ยังไม่ฝังคะแนน}

  \tccase{1}{ส่งแบบประเมินครบแล้วเห็นผล}{%
    \begin{tcsteps}
      \item เข้าเมนู ``ประเมิน'' จากแดชบอร์ดนิสิต
      \item อ่านคำอธิบาย/เกณฑ์ก่อนเริ่มทำ
      \item ตอบคำถามทุกข้อแล้วกด ``ส่งคำตอบ''
    \end{tcsteps}}{%
    บันทึกสำเร็จ และ JSON/หน้าจอมีคะแนนรวมกับป้ายแปลผล}{%
    หน่วยทดสอบ \tccode{buildStudentFacingSurveyResult} ผ่าน และ API ส่งสรุปผลหลัง submit}{Pass}

  \tccase{2}{ตัวเลือกที่นิสิตเห็นไม่มีคะแนนฝังใน options}{%
    \begin{tcsteps}
      \item เรียก \tccode{GET /api/student/surveys/active} ด้วยเซสชันนิสิต
      \item ตรวจฟิลด์ \tccode{questions[].options}
    \end{tcsteps}}{%
    แต่ละ option มีเพียง label และ value ไม่มีฟิลด์ score}{%
    แต่ละ option มีเพียง label และ value ไม่มีฟิลด์ score}{Pass}

  \tccase{3}{หลังส่งแล้วยังเห็นเกณฑ์และผลซ้ำได้}{%
    \begin{tcsteps}
      \item เปิด ``ดูผลและเกณฑ์อีกครั้ง''
      \item ตรวจว่าคำอธิบายเกณฑ์และสรุปคะแนนยังแสดง
    \end{tcsteps}}{%
    เห็นเกณฑ์ + คะแนน/แปลผล โดยไม่ต้องทำแบบใหม่}{%
    UI หน้า assessment และ modal รองรับ}{Pass}
\end{testscenario}

% TC03 — บทบาทเพิ่มเติมและการเข้าถึงนิสิต (IDOR)
\begin{testscenario}{3}{3}
  \tcmeta{TC03}{CS 2568/01}{ระบบเว็บแอปพลิเคชัน ``NUCARE''}{นายสมชาย ใจดี}{วันที่ 22 ตุลาคม พ.ศ.\ 2569}{1.0}
  \tcname{การตรวจสอบสิทธิ์บทบาทเพิ่มเติมและการเข้าถึงข้อมูลนิสิต}
  \tcpre{มีหน่วยทดสอบ \tccode{roles.test.ts} และ \tccode{access-control.test.ts}
    ในเซิร์ฟเวอร์ และรันได้ด้วย \tccode{npm run test:unit}}
  \tcdesc{ทดสอบว่าการตรวจบทบาทรองรับทั้งบทบาทหลักและบทบาทเพิ่มเติม
    และการเข้าถึงนิสิตรายคนถูกจำกัดตามการมอบหมาย (ป้องกัน IDOR) สำหรับที่ปรึกษา/ผู้สอน}

  \tccase{1}{ตรวจสอบ hasAssignedRole เมื่อมีบทบาทเพิ่มเติมใน payload}{%
    \begin{tcsteps}
      \item เปิดไฟล์ \tccode{src/lib/roles.test.ts}
      \item รันกรณี RL-02 ด้วยผู้ใช้ role=ADVISOR และ roles รวม ADMIN
      \item ตรวจว่า \tccode{hasAssignedRole} สำหรับ ADMIN เป็นจริง
    \end{tcsteps}}{%
    ระบบนับทั้งบทบาทหลักและรายการ roles แล้วคืนค่า true สำหรับ ADMIN}{%
    รัน test:unit ผ่านกรณี RL-02 ตามที่คาดหวัง}{Pass}

  \tccase{2}{ตรวจสอบ canAccessStudent ปฏิเสธที่ปรึกษาที่ไม่มี assignment}{%
    \begin{tcsteps}
      \item เปิดไฟล์ \tccode{src/lib/access-control.test.ts}
      \item ตั้ง mock ให้ findActiveAdvisorAssignment คืนค่า false
      \item เรียก canAccessStudent ของที่ปรึกษาที่ไม่มี assignment
    \end{tcsteps}}{%
    ฟังก์ชันคืนค่า false และไม่อนุญาตเข้าถึงนิสิตที่ไม่อยู่ในความดูแล}{%
    รัน test:unit ผ่านกรณี AC-03 ตามที่คาดหวัง}{Pass}

  \tccase{3}{ตรวจสอบผู้บริหารไม่เข้าถึงนิสิตรายคนผ่าน canAccessStudent}{%
    \begin{tcsteps}
      \item เรียก canAccessStudent ด้วยบทบาท EXECUTIVE ในหน่วยทดสอบ AC-05
      \item ตรวจค่าที่คืน
    \end{tcsteps}}{%
    ฟังก์ชันคืนค่า false ตามนโยบาย PDPA / ไม่เปิดประวัติรายบุคคล}{%
    รัน test:unit ผ่านกรณี AC-05 ตามที่คาดหวัง}{Pass}
\end{testscenario}

% TC04 — นัดหมายสุขภาพจิต (persist)
\begin{testscenario}{4}{4}
  \tcmeta{TC04}{CS 2568/01}{ระบบเว็บแอปพลิเคชัน ``NUCARE''}{นายสมชาย ใจดี}{วันที่ 22 ตุลาคม พ.ศ.\ 2569}{1.0}
  \tcname{การบันทึกนัดหมายของทีมสุขภาพจิตและแสดงให้นิสิต}
  \tcpre{มีเคส referral สถานะ ACCEPTED ไปยังศูนย์ให้คำปรึกษา และรัน migration MentalHealthAppointment แล้ว}
  \tcdesc{ทดสอบการสร้างนัดหมายที่บันทึกลงฐานข้อมูล การแปลงวันเวลา และหน่วยทดสอบตัวช่วย serialize}

  \tccase{1}{ตรวจสอบ buildScheduledAt / serialize ในหน่วยทดสอบ}{%
    \begin{tcsteps}
      \item รัน \tccode{npm run test:unit} ในโฟลเดอร์ server
      \item ตรวจกรณี MHA-01 ถึง MHA-04
    \end{tcsteps}}{%
    ทุกกรณีผ่าน และ DTO มี date/time/location/providerName ครบ}{%
    รัน test:unit ผ่านกรณี MHA ทั้งหมด}{Pass}

  \tccase{2}{สร้างนัดหมายจากหน้านักจิตแล้วโหลดกลับจาก dashboard}{%
    \begin{tcsteps}
      \item เข้าสู่ระบบด้วยบทบาท MENTAL\_HEALTH
      \item เปิดเคสที่รับแล้ว แล้วสร้างนัดหมายวันเวลา/สถานที่
      \item รีเฟรชหน้า dashboard
    \end{tcsteps}}{%
    นัดหมายยังอยู่หลังรีเฟรช และปรากฏในรายการนัด}{%
    API สร้างนัดหมายและ GET dashboard คืน appointments}{Pass}
\end{testscenario}

% TC05 — รายงานและตั้งค่าแอดมิน
\begin{testscenario}{5}{5}
  \tcmeta{TC05}{CS 2568/01}{ระบบเว็บแอปพลิเคชัน ``NUCARE''}{นายสมชาย ใจดี}{วันที่ 22 ตุลาคม พ.ศ.\ 2569}{1.0}
  \tcname{การดูสรุปรายงานและทางลัดตั้งค่าของผู้ดูแลระบบ}
  \tcpre{มีบัญชี ADMIN และเซิร์ฟเวอร์พร้อมข้อมูลนิสิต/แบบประเมินอย่างน้อยบางส่วน}
  \tcdesc{ทดสอบว่าหน้า reports ดึงสรุปแบบไม่ระบุตัวตนได้ และหน้า settings ลิงก์ไปเครื่องมือจริง}

  \tccase{1}{เปิดหน้าสถิติและรายงาน}{%
    \begin{tcsteps}
      \item เข้าสู่ระบบด้วย ADMIN
      \item เปิด \tccode{/admin?view=reports}
      \item กดรีเฟรชสรุป
    \end{tcsteps}}{%
    แสดงจำนวนนิสิต คัดกรองแล้ว referral/SOS ที่เปิดอยู่ โดยไม่มีชื่อนิสิต}{%
    API reports/summary คืน overall และ operations}{Pass}

  \tccase{2}{เปิดหน้าตั้งค่าและเข้าแบบประเมิน}{%
    \begin{tcsteps}
      \item เปิด \tccode{/admin?view=settings}
      \item กดการ์ดแบบประเมินและรอบคัดกรอง
    \end{tcsteps}}{%
    ระบบพาไป \tccode{/admin/surveys} ได้}{%
    ลิงก์ใน SettingsView ทำงานตามที่คาดหวัง}{Pass}
\end{testscenario}

% TC06 — สวัสดิการพัฒนานิสิต
\begin{testscenario}{6}{3}
  \tcmeta{TC06}{CS 2568/01}{ระบบเว็บแอปพลิเคชัน ``NUCARE''}{นายสมชาย ใจดี}{วันที่ 6 พฤศจิกายน พ.ศ.\ 2569}{1.1}
  \tcname{การอัปเดตสถานะ referral ของงานพัฒนานิสิต}
  \tcpre{มี CaseReferral ที่ target = STUDENT\_DEVELOPMENT สถานะ PENDING และบัญชี STUDENT\_DEVELOPMENT พร้อมใช้งาน}
  \tcdesc{ทดสอบการรับงาน ปิดงาน ตามลำดับสถานะ พร้อม ownership และแจ้งอาจารย์ที่ส่งต่อ}

  \tccase{1}{รับงาน referral จากหน้าสวัสดิการ}{%
    \begin{tcsteps}
      \item เข้าสู่ระบบด้วย STUDENT\_DEVELOPMENT
      \item เปิดมุมมองสวัสดิการ
      \item กด ``รับงาน'' บนเคส PENDING แล้วยืนยัน
    \end{tcsteps}}{%
    สถานะเป็น ACCEPTED / กำลังประสานงาน บันทึก acceptedByUserId และแจ้งอาจารย์}{%
    อัปเดตสถานะ referral เป็น ACCEPTED สำเร็จ พร้อม toast รับงานสำเร็จ}{Pass}

  \tccase{2}{ปิดงานโดยผู้รับเคส}{%
    \begin{tcsteps}
      \item จากเคส ACCEPTED ที่ตนรับไว้
      \item กด ``ปิดงาน'' แล้วยืนยัน
    \end{tcsteps}}{%
    สถานะเป็น CLOSED / ปิดงานแล้ว และแจ้งอาจารย์}{%
    ปิดงานสำเร็จ สถานะ CLOSED}{Pass}

  \tccase{3}{ปฏิเสธการเปลี่ยนสถานะผิดลำดับ}{%
    \begin{tcsteps}
      \item พยายามอัปเดตเคส CLOSED กลับเป็น ACCEPTED ผ่าน API
    \end{tcsteps}}{%
    API ตอบ 400 พร้อมข้อความว่าเปลี่ยนสถานะไม่ได้}{%
    ได้ 400 ตาม transition guard}{Pass}
\end{testscenario}

% TC07 — สัญญาณจากผู้สอน / รับเรื่องติดตาม
\begin{testscenario}{7}{7}
  \tcmeta{TC07}{CS 2568/01}{ระบบเว็บแอปพลิเคชัน ``NUCARE''}{นายสมชาย ใจดี}{วันที่ 28 ตุลาคม พ.ศ.\ 2569}{1.1}
  \tcname{การรับและติดตามข้อสังเกตจากผู้สอนโดยอาจารย์ที่ปรึกษา}
  \tcpre{มีนิสิตในความดูแล มี BehaviorRecord จากผู้สอน และ Alert ประเภท BEHAVIOR\_REPORTED ถึงที่ปรึกษา}
  \tcdesc{ทดสอบว่าที่ปรึกษาเห็นคิวจากผู้สอน กดรับเรื่อง/ปิดเรื่องได้ และสถานะสะท้อนฝั่งผู้สอน}

  \tccase{1}{เปิดหน้าสัญญาณจากผู้สอน}{%
    \begin{tcsteps}
      \item เข้าสู่ระบบด้วย ADVISOR
      \item เปิดเมนู ``สัญญาณจากผู้สอน''
      \item ตรวจคิวรับเรื่องและรายชื่อนิสิตที่มีบันทึก
    \end{tcsteps}}{%
    แสดงคิว Alert จากผู้สอนและรายชื่อที่มี behaviorRecordCount $>$ 0}{%
    AcademicSignalsView แสดงคิวและรายชื่อตามข้อมูล API}{Pass}

  \tccase{2}{รับเรื่องแล้วปิดเรื่อง}{%
    \begin{tcsteps}
      \item กด ``รับเรื่อง'' บน Alert ที่สถานะ OPEN
      \item กด ``ปิดเรื่อง'' หลังติดตาม
      \item เปิดประวัติฝั่งผู้สอนของรายการเดิม
    \end{tcsteps}}{%
    Alert เป็น ACKNOWLEDGED แล้ว CLOSED; ฝั่งผู้สอนเห็นสถานะรับเรื่องแล้ว / ดำเนินการเสร็จสิ้น}{%
    \tccode{PATCH /api/alerts/:id/acknowledge} และ \tccode{/close} อัปเดตได้เฉพาะผู้รับ; follow-up status สะท้อนถูกต้อง}{Pass}
\end{testscenario}

% TC08 — แชทเพื่อน AI ของนิสิต
\begin{testscenario}{8}{8}
  \tcmeta{TC08}{CS 2568/01}{ระบบเว็บแอปพลิเคชัน ``NUCARE''}{นายสมชาย ใจดี}{วันที่ 28 ตุลาคม พ.ศ.\ 2569}{1.1}
  \tcname{การคุยกับเพื่อน AI และการกันถามการบ้าน}
  \tcpre{มีบัญชีนิสิตที่เข้าสู่ระบบได้ และเซิร์ฟเวอร์ตั้ง \tccode{CHAT_PROVIDER=mock}}
  \tcdesc{ทดสอบการเลือกการ์ดพี่ต้น/พี่เมย์ การแชทปกติ การปฏิเสธคำขอทำการบ้าน และประกาศความเป็นส่วนตัวว่าไม่เก็บข้อความ}

  \tccase{1}{เลือก persona และเริ่มคุย}{%
    \begin{tcsteps}
      \item เข้าสู่ระบบด้วยบทบาท STUDENT
      \item เปิด \tccode{/student/chat}
      \item กดการ์ดพี่ต้นหรือพี่เมย์เพื่อเริ่มคุย
    \end{tcsteps}}{%
    แสดงข้อความทักทายของ persona ที่เลือก}{%
    หน้าแชทโหลด persona และแสดง greeting}{Pass}

  \tccase{2}{ระบบปฏิเสธคำขอทำการบ้าน}{%
    \begin{tcsteps}
      \item ในห้องแชท พิมพ์ข้อความขอช่วยทำการบ้านหรือขอเฉลยข้อสอบ
      \item กดส่ง
    \end{tcsteps}}{%
    ได้ข้อความปฏิเสธสุภาพ และไม่ให้เฉลย/โค้ด}{%
    หน่วยทดสอบ CG-01/MP-01 คืน blockedReason=homework}{Pass}

  \tccase{3}{คุยเรื่องเครียดได้ตามปกติ}{%
    \begin{tcsteps}
      \item ส่งข้อความว่าเครียดหรือเหนื่อย
    \end{tcsteps}}{%
    ได้คำตอบรับฟังโทนเพื่อน ไม่ถูกบล็อก}{%
    หน่วยทดสอบ MP-02 ผ่าน}{Pass}
\end{testscenario}

% TC09 — สัญญาณแนวโน้มสุขภาวะจากแชท (ไม่เก็บข้อความ)
\begin{testscenario}{9}{4}
  \tcmeta{TC09}{CS 2568/01}{ระบบเว็บแอปพลิเคชัน ``NUCARE''}{นายสมชาย ใจดี}{วันที่ 28 ตุลาคม พ.ศ.\ 2569}{1.0}
  \tcname{Wellness Trend Indicators จากแชทเพื่อน AI}
  \tcpre{มีหน่วยทดสอบ \tccode{wellness-signals.test.ts} และ schema สรุป/จุดตรวจจับแนวโน้ม}
  \tcdesc{ตรวจว่าสกัดรหัสแนวโน้มจากข้อความได้โดยไม่คืนหรือเก็บเนื้อหา และรองรับหลายสัญญาณในข้อความเดียว}

  \tccase{1}{ตรวจจับประเด็นอารมณ์ซึมเศร้า}{%
    \begin{tcsteps}
      \item เปิดไฟล์ \tccode{src/lib/student-chat/wellness-signals.ts}
      \item รันกรณี WT-01 ด้วยข้อความที่มีคำว่าซึมเศร้า/ท้อแท้
    \end{tcsteps}}{%
    คืนรหัส \tccode{DEPRESSIVE_THEME} และมีชื่อทางการภาษาไทย}{%
    หน่วยทดสอบ WT-01 ผ่าน}{Pass}

  \tccase{2}{หลายแนวโน้มในข้อความเดียว}{%
    \begin{tcsteps}
      \item รันกรณี WT-02 ด้วยข้อความเครียดสอบ + นอนไม่หลับ + เหงา
    \end{tcsteps}}{%
    คืน \tccode{ACADEMIC\_STRESS}, \tccode{SLEEP\_DISTURBANCE}, \tccode{LONELINESS\_ISOLATION}}{%
    หน่วยทดสอบ WT-02 ผ่าน}{Pass}

  \tccase{3}{สัญญาณวิกฤต}{%
    \begin{tcsteps}
      \item รันกรณี WT-03 ด้วยภาษาอยากทำร้ายตนเอง
    \end{tcsteps}}{%
    มีรหัส \tccode{SELF\_HARM\_CRISIS} ในผลลัพธ์}{%
    หน่วยทดสอบ WT-03 ผ่าน}{Pass}

  \tccase{4}{ข้อความกลาง ๆ ไม่สร้างสัญญาณ}{%
    \begin{tcsteps}
      \item รันกรณี WT-04 และ WT-05
    \end{tcsteps}}{%
    คืนอาร์เรย์ว่าง}{%
    หน่วยทดสอบ WT-04/WT-05 ผ่าน}{Pass}
\end{testscenario}

% TC10 — คลื่นแบบประเมิน + กฎกลุ่มเป้าหมายแบบไม่เปิดเผยเหตุผล
\begin{testscenario}{10}{3}
  \tcmeta{TC10}{CS 2568/01}{ระบบเว็บแอปพลิเคชัน ``NUCARE''}{นายสมชาย ใจดี}{วันที่ 28 ตุลาคม พ.ศ.\ 2568}{1.0}
  \tcname{การตั้งค่าคลื่นแบบประเมินและกฎคัดเลือกกลุ่มเป้าหมายโดยไม่เปิดเผยเหตุผลให้นิสิต}
  \tcpre{มีบัญชีผู้ดูแลระบบ และมีข้อมูลนิสิตที่มีผลเสี่ยง/สัญญาณแชทต่างกัน}
  \tcdesc{ทดสอบว่าแอดมินตั้งความถี่ กลุ่มเป้าหมาย และเปิดคลื่นได้
    และนิสิตที่ไม่เข้าเกณฑ์จะไม่เห็นแบบประเมินหรือข้อความบอกเหตุผลคัดเลือก}

  \tccase{1}{กฎ RULE\_BASED คัดกลุ่มตามความเสี่ยงและสัญญาณแชท}{%
    \begin{tcsteps}
      \item รัน unit ของ \tccode{survey-eligibility}
      \item ตรวจเคส ALL\_STUDENTS ผ่านทุกคน และ RULE\_BASED ว่างแล้วปฏิเสธ
      \item ตรวจเคส prior risk + wellness OR ภายในมิติ / AND ข้ามมิติ
    \end{tcsteps}}{%
    เคส unit ผ่านครบตาม assertion}{%
    เคส unit ผ่านครบตาม assertion}{Pass}

  \tccase{2}{แอดมินดูจำนวนที่เข้าเกณฑ์โดยไม่ได้รับรายชื่อใน preview}{%
    \begin{tcsteps}
      \item เข้าหน้าจัดการแบบสอบถาม ตั้ง audience เป็นเฉพาะกลุ่มตามเงื่อนไข
      \item กด ``ดูจำนวนที่เข้าเกณฑ์''
    \end{tcsteps}}{%
    API คืนเฉพาะ \tccode{matchedCount} ไม่คืนรายชื่อนิสิต}{%
    API คืนเฉพาะ \tccode{matchedCount} ไม่คืนรายชื่อนิสิต}{Pass}

  \tccase{3}{นิสิตที่ไม่เข้าเกณฑ์ไม่เห็นแบบและไม่ถูกบอกเหตุผล}{%
    \begin{tcsteps}
      \item เปิดคลื่นแบบ RULE\_BASED ที่กรองเฉพาะความเสี่ยงสูง
      \item เข้าสู่ระบบด้วยนิสิตที่ความเสี่ยงปกติและไม่มีสัญญาณตรงเงื่อนไข
      \item เรียก \tccode{GET /api/student/surveys/active}
    \end{tcsteps}}{%
    ได้รายการว่างหรือไม่มีแบบประเมิน และไม่มีข้อความว่าไม่เข้ากลุ่มเสี่ยง}{%
    ได้รายการว่างหรือไม่มีแบบประเมิน และไม่มีข้อความว่าไม่เข้ากลุ่มเสี่ยง}{Pass}
\end{testscenario}

% TC11 — ตัดคำนำหน้าชื่อก่อนย่อตัวอักษร + อวตาร์ไม่ใช้ท้ายรหัส
\begin{testscenario}{11}{3}
  \tcmeta{TC11}{CS 2568/01}{ระบบเว็บแอปพลิเคชัน ``NUCARE''}{นายสมชาย ใจดี}{วันที่ 8 พฤศจิกายน พ.ศ.\ 2569}{1.1}
  \tcname{การย่อชื่อและอวตาร์รายชื่อโดยไม่ใช้ท้ายรหัสนิสิต}
  \tcpre{มีฟังก์ชัน \tccode{getNameInitials} / \tccode{resolveRosterAvatarMark} ในเซิร์ฟเวอร์}
  \tcdesc{ทดสอบว่าการย่อชื่อใช้เฉพาะตัวอักษรแรกของชื่อและนามสกุล
    ข้ามคำนำหน้า และอวตาร์รายชื่อไม่ดึงท้ายรหัสนิสิต เช่น \tccode{-0}}

  \tccase{1}{ตัดคำนำหน้านักศึกษาและบุคลากร}{%
    \begin{tcsteps}
      \item รัน unit ของ \tccode{name-prefix}
      \item ตรวจชื่อเช่น นายสมชาย ใจดี และ ผศ.ดร.วิชัย ใจดี
    \end{tcsteps}}{%
    ได้ตัวย่อ สจ และ วจ ตามลำดับ ไม่มีตัวอักษรจากคำนำหน้า}{%
    unit ผ่าน}{Pass}

  \tccase{2}{ข้ามสระนำไทยตอนเลือกพยัญชนะ}{%
    \begin{tcsteps}
      \item ตรวจชื่อที่มีนามสกุลขึ้นต้นด้วยสระนำ เช่น ใจดี
    \end{tcsteps}}{%
    ย่อด้วยพยัญชนะ จ ไม่ใช่สระ ใ}{%
    unit ผ่าน}{Pass}

  \tccase{3}{อวตาร์ไม่ใช้ท้ายรหัสนิสิต}{%
    \begin{tcsteps}
      \item เรียก \tccode{resolveRosterAvatarMark} ด้วยชื่อนิสิตและรหัสที่ลงท้าย \tccode{-0}
    \end{tcsteps}}{%
    คืนตัวย่อจากชื่อ ไม่ใช่ \tccode{-0}}{%
    unit ผ่าน}{Pass}
\end{testscenario}

% TC12 — สรุปสถานะติดตาม BehaviorRecord จาก Alert
\begin{testscenario}{12}{12}
  \tcmeta{TC12}{CS 2568/01}{ระบบเว็บแอปพลิเคชัน ``NUCARE''}{นายสมชาย ใจดี}{วันที่ 29 ตุลาคม พ.ศ.\ 2569}{1.1}
  \tcname{การสรุปสถานะติดตามข้อสังเกตจากสถานะ Alert ของที่ปรึกษา}
  \tcpre{รัน unit test ของ \tccode{getBehaviorRecordFollowUpStatus}}
  \tcdesc{ทดสอบว่าสถานะ PENDING / ACKNOWLEDGED / COMPLETED สอดคล้องกับชุด Alert เมื่อมีหลายที่ปรึกษา}

  \tccase{1}{ไม่มี Alert หรือยัง OPEN ทั้งหมด}{%
    \begin{tcsteps}
      \item เรียก \tccode{getBehaviorRecordFollowUpStatus([])}
      \item เรียกด้วยชุด OPEN อย่างเดียว
    \end{tcsteps}}{คืน PENDING}{unit ผ่าน}{Pass}

  \tccase{2}{มี ACKNOWLEDGED หรือ CLOSED บางใบ}{%
    \begin{tcsteps}
      \item เรียกด้วย OPEN + ACKNOWLEDGED
      \item เรียกด้วย OPEN + CLOSED
    \end{tcsteps}}{คืน ACKNOWLEDGED}{unit ผ่าน}{Pass}

  \tccase{3}{ปิดครบทุกใบ}{%
    \begin{tcsteps}
      \item เรียกด้วย CLOSED + CLOSED
    \end{tcsteps}}{คืน COMPLETED}{unit ผ่าน}{Pass}
\end{testscenario}

% TC13 — ยุบบทบาทผู้สอนเข้าที่ปรึกษา + กฎแจ้งข้อสังเกตนอกสังกัด
\begin{testscenario}{13}{4}
  \tcmeta{TC13}{CS 2568/01}{ระบบเว็บแอปพลิเคชัน ``NUCARE''}{นายสมชาย ใจดี}{วันที่ 8 พฤศจิกายน พ.ศ.\ 2569}{1.1}
  \tcname{การยุบบทบาทผู้สอนและการแจ้งข้อสังเกตเฉพาะนิสิตนอกสังกัด}
  \tcpre{มีบัญชี ADVISOR และนิสิตในระบบ; unit test roles/access-control และ \tccode{behavior-observation.test.ts} ผ่าน}
  \tcdesc{ตรวจว่า INSTRUCTOR ถูกยุบเป็น ADVISOR · แจ้งข้อสังเกตได้เฉพาะนิสิตที่ไม่ใช่ในสังกัดผู้แจ้ง · สัญญาณเข้าที่ปรึกษาของนิสิต}

  \tccase{1}{ยุบบทบาทใน getAssignedRoles}{%
    \begin{tcsteps}
      \item เรียก \tccode{getAssignedRoles} ด้วย role = INSTRUCTOR
      \item ตรวจ \tccode{hasAssignedRole(INSTRUCTOR, ADVISOR)}
    \end{tcsteps}}{ได้บทบาท ADVISOR และผ่านการตรวจสิทธิ์ที่ปรึกษา}{unit RL-08 ผ่าน}{Pass}

  \tccase{2}{ที่ปรึกษาไม่มีนิสิตในดูแล}{%
    \begin{tcsteps}
      \item เข้าสู่ระบบด้วย \tccode{advisor.peer@nucare.local} (ก่อน walkthrough seed)
      \item ตรวจเมนูด้านซ้าย
    \end{tcsteps}}{ไม่แสดงเมนู ``นักศึกษาในที่ปรึกษา'' แต่มี ``แจ้งข้อสังเกต''}{UI ซ่อนเมื่อ students.length = 0}{Pass}

  \tccase{3}{แจ้งข้อสังเกตนอกสังกัดแล้วเข้าสัญญาณ}{%
    \begin{tcsteps}
      \item เปิดเมนูแจ้งข้อสังเกต
      \item ค้นหา — รายชื่อไม่รวมนิสิตในสังกัดตัวเอง
      \item เลือกนิสิตของอาจารย์ท่านอื่นแล้วส่ง
    \end{tcsteps}}{สร้าง BehaviorRecord และ Alert ชนิด BEHAVIOR\_REPORTED ถึงที่ปรึกษาของนิสิต (เมนูสัญญาณเข้า)}{API \tccode{POST /api/advisor/observation-records} ปฏิเสธ own advisee}{Pass}

  \tccase{4}{ห้ามแจ้งนิสิตในสังกัดตัวเอง}{%
    \begin{tcsteps}
      \item รัน \tccode{behavior-observation.test.ts} เคส BO-01
      \item POST observation ด้วย studentId ที่เป็น advisee ของผู้แจ้ง
    \end{tcsteps}}{\tccode{canReportObservationAboutStudent(true)=false} · API ได้ 400 \tccode{OWN\_ADVISEE\_NOT\_ALLOWED}}{ผ่าน}{Pass}
\end{testscenario}

% TC14 — roadmap ข้อสังเกตจากเหตุการณ์จริงเท่านั้น
\begin{testscenario}{14}{15}
  \tcmeta{TC14}{CS 2568/01}{ระบบเว็บแอปพลิเคชัน ``NUCARE''}{นายสมชาย ใจดี}{วันที่ 2 พฤศจิกายน พ.ศ.\ 2569}{1.3}
  \tcname{การสร้างเส้นทางติดตามข้อสังเกตจากเหตุการณ์ที่มีจริง}
  \tcpre{รัน unit test ของ \tccode{buildObservationRoadmap}}
  \tcdesc{ทดสอบว่า roadmap เริ่มจากขั้นแจ้งเท่านั้น เพิ่มรับเรื่อง/ปิดเรื่องพร้อมเหตุผล ส่งต่อผู้เชี่ยวชาญต่อได้หลังปิด และรวมหมุดส่งต่อปลายทางเดียวกัน}

  \tccase{1}{ยังไม่มีใครรับเรื่อง}{%
    \begin{tcsteps}
      \item เรียก \tccode{buildObservationRoadmap} ด้วย Alert สถานะ OPEN
      \item ตรวจจำนวนและป้ายขั้น
    \end{tcsteps}}{มีเพียงขั้น ``แจ้งข้อสังเกต''}{unit ผ่าน}{Pass}

  \tccase{2}{ที่ปรึกษารับเรื่องบางใบ}{%
    \begin{tcsteps}
      \item เรียกด้วย ACKNOWLEDGED ร่วมกับ OPEN
      \item ตรวจลำดับป้ายและชื่อผู้รับเรื่อง
    \end{tcsteps}}{ได้ขั้นแจ้งแล้วตามด้วยรับเรื่องเฉพาะใบที่ดำเนินการ}{unit ผ่าน}{Pass}

  \tccase{3}{ปิดเรื่องเก็บเหตุผลและยังส่งต่อได้}{%
    \begin{tcsteps}
      \item เรียกด้วย CLOSED พร้อม \tccode{resolutionNote}
      \item เพิ่ม referral หลังเวลาปิดเรื่อง
      \item ตรวจลำดับ kind และรายละเอียดเหตุผล
    \end{tcsteps}}{ได้แจ้ง $\rightarrow$ รับ $\rightarrow$ ปิด (มีเหตุผล) $\rightarrow$ ส่งต่อสองบรรทัด}{unit ผ่าน}{Pass}

  \tccase{4}{รวมส่งต่อปลายทางซ้ำเป็นหมุดเดียว}{%
    \begin{tcsteps}
      \item เรียกด้วย referral ปลายทางเดียวกัน 2 รายการ
      \item ตรวจจำนวนหมุด kind referred และข้อความจำนวนครั้ง
    \end{tcsteps}}{มีหมุดส่งต่อเพียง 1 อัน แสดงสถานะล่าสุดและจำนวนครั้ง}{unit ผ่าน}{Pass}

  \tccase{5}{UI รวมหมุดซ้ำจากป้ายเก่าและตัดสองบรรทัด}{%
    \begin{tcsteps}
      \item เรียก \tccode{prepareRoadmapSteps} ด้วยป้าย ``ส่งต่อ · \ldots'' ซ้ำ
      \item ตรวจป้ายสองบรรทัดและจำนวนหมุด
    \end{tcsteps}}{เหลือหมุดส่งต่อ 1 อัน ป้ายขึ้นบรรทัดใหม่}{unit ผ่าน}{Pass}
\end{testscenario}

% TC15 — ความถี่แบบสอบถามแบบกำหนดเองและร่างคลื่นรอบถัดไป
\begin{testscenario}{15}{15}
  \tcmeta{TC15}{CS 2568/01}{ระบบเว็บแอปพลิเคชัน ``NUCARE''}{นายสมชาย ใจดี}{วันที่ 29 ตุลาคม พ.ศ.\ 2569}{1.0}
  \tcname{การตั้งความถี่แบบสอบถามแบบกำหนดเองและการสร้างร่างคลื่นรอบถัดไป}
  \tcpre{รัน unit test ของ \tccode{survey-schedule}}
  \tcdesc{ทดสอบว่าช่วงกำหนดเองแปลงเป็นวันได้ถูกต้อง และร่างคลื่นไม่บังคับเปิดอัตโนมัติล่วงหน้าโดยไม่มีข้อมูล}

  \tccase{1}{แปลงหน่วยวัน สัปดาห์ เดือน}{%
    \begin{tcsteps}
      \item เรียก \tccode{scheduleIntervalToDays}
      \item ตรวจค่า 3 วัน 2 สัปดาห์ และ 1 เดือน
    \end{tcsteps}}{ได้ 3, 14 และ 30 ตามลำดับ}{unit ผ่าน}{Pass}

  \tccase{2}{CUSTOM\_INTERVAL ต้องมีจำนวนและหน่วย}{%
    \begin{tcsteps}
      \item เรียก \tccode{validateCustomSchedule} โดยขาดค่า
    \end{tcsteps}}{คืนข้อความผิดพลาดที่อ่านรู้เรื่อง}{unit ผ่าน}{Pass}

  \tccase{3}{ร่างคลื่นตามเทอมใช้ termKey ที่ใส่}{%
    \begin{tcsteps}
      \item เรียก \tccode{buildNextWaveDraft} ด้วย ACADEMIC\_TERM และ termKey
    \end{tcsteps}}{label และ termKey ตรงค่าที่กำหนด}{unit ผ่าน}{Pass}
\end{testscenario}

% TC16 — ผู้ให้บริการอีเมลและช่องทางแจ้งเตือนผู้ใช้
\begin{testscenario}{16}{16}
  \tcmeta{TC16}{CS 2568/01}{ระบบเว็บแอปพลิเคชัน ``NUCARE''}{นายสมชาย ใจดี}{วันที่ 29 ตุลาคม พ.ศ.\ 2569}{1.0}
  \tcname{โครงส่งอีเมลแบบเปลี่ยน provider ได้และการเลือกช่องทางแจ้งเตือน}
  \tcpre{รัน unit test ของ \tccode{email} และ \tccode{notify-user}}
  \tcdesc{ทดสอบว่าเลือก provider จาก env ได้ ข้ามอีเมลที่ผิดรูปแบบ และส่ง in-app/email ตามช่องทางที่เลือก}

  \tccase{1}{ค่าเริ่มต้นเป็น console provider}{%
    \begin{tcsteps}
      \item ไม่ตั้ง \tccode{EMAIL\_PROVIDER}
      \item เรียก \tccode{resolveEmailProviderKind}
    \end{tcsteps}}{ได้ค่า \tccode{console}}{unit ผ่าน}{Pass}

  \tccase{2}{ข้ามอีเมลผู้รับที่ไม่ถูกต้อง}{%
    \begin{tcsteps}
      \item เรียก \tccode{sendAppEmail} ด้วยที่อยู่ไม่มี @
    \end{tcsteps}}{\tccode{ok=false} และ \tccode{skipped=true}}{unit ผ่าน}{Pass}

  \tccase{3}{notifyUser ส่งเมลเมื่อเปิดช่องทาง email}{%
    \begin{tcsteps}
      \item ตั้ง \tccode{EMAIL\_PROVIDER=console}
      \item เรียก \tccode{notifyUser} ด้วย \tccode{channels.email=true}
    \end{tcsteps}}{\tccode{emailSent=true}}{unit ผ่าน}{Pass}
\end{testscenario}

% TC17 — สัญญาณแชท AI และการมองเห็นผลละเอียดอ่อนตามบทบาท
\begin{testscenario}{17}{17}
  \tcmeta{TC17}{CS 2568/01}{ระบบเว็บแอปพลิเคชัน ``NUCARE''}{นายสมชาย ใจดี}{วันที่ 29 ตุลาคม พ.ศ.\ 2569}{1.0}
  \tcname{สัญญาณจากแชท AI ในประวัตินิสิตและสิทธิ์มองเห็นผลละเอียดอ่อน}
  \tcpre{รัน unit test ของ \tccode{access-control}, \tccode{notify-advisors-of-wellness}, \tccode{student-wellness-summary}}
  \tcdesc{ทดสอบว่าผู้บริหารเข้าถึงนิสิตได้ทั้งหมด สัญญาณแชทไม่สร้างแจ้งเตือนในกล่องกระดิ่ง และไม่รั่วเนื้อหาแชท}

  \tccase{1}{EXECUTIVE เข้าถึงนิสิตรายคนได้}{%
    \begin{tcsteps}
      \item เรียก \tccode{canAccessStudent} ด้วยบทบาท EXECUTIVE
    \end{tcsteps}}{\tccode{true}}{unit ผ่าน}{Pass}

  \tccase{2}{สัญญาณแชทไม่สร้าง Alert ในกล่องแจ้งเตือน}{%
    \begin{tcsteps}
      \item เรียก \tccode{notifyAdvisorsOfWellnessSignals} ด้วย \tccode{SELF\_HARM\_CRISIS}
      \item ตรวจว่าไม่เรียก \tccode{alert.createMany}
    \end{tcsteps}}{คืน 0 และไม่สร้าง Alert — สัญญาณแสดงในหน้ารายละเอียดนิสิตเท่านั้น}{unit ผ่าน}{Pass}

  \tccase{3}{STUDENT\_DEVELOPMENT ไม่ได้ scope นิสิตรายคนผ่าน access-control}{%
    \begin{tcsteps}
      \item เรียก \tccode{getStudentScopeWhere} ด้วย STUDENT\_DEVELOPMENT
    \end{tcsteps}}{ได้ sentinel \tccode{\_\_no\_student\_scope\_\_}}{unit ผ่าน}{Pass}
\end{testscenario}

% TC18 — ปิดแจ้งเตือนสัญญาณโดยไม่ถือว่ารับเรื่อง / บังคับรับเรื่องก่อนปิดเคส
\begin{testscenario}{18}{18}
  \tcmeta{TC18}{CS 2568/01}{ระบบเว็บแอปพลิเคชัน ``NUCARE''}{นายสมชาย ใจดี}{วันที่ 8 พฤศจิกายน พ.ศ.\ 2569}{1.1}
  \tcname{นโยบายปิดแจ้งเตือน: dismiss ไม่ถือว่ารับเรื่อง และ follow\_up ต้องรับเรื่องก่อน}
  \tcpre{รัน unit test ของ \tccode{alert-close-policy}}
  \tcdesc{แยกโหมดปิดกระดิ่ง/สัญญาณ ออกจากการปิดเรื่องหลังรับติดตาม และบังคับลำดับรับเรื่องก่อนบันทึกและปิดเคส}

  \tccase{1}{dismiss\_signal ไม่ตั้ง acknowledgedAt}{%
    \begin{tcsteps}
      \item เรียก \tccode{resolveAlertCloseData} ด้วย \tccode{mode=dismiss\_signal} และยังไม่เคยรับเรื่อง
      \item ตรวจฟิลด์ที่คืน
    \end{tcsteps}}{allowed พร้อม \tccode{closedAt} แต่ไม่มี \tccode{acknowledgedAt} และ audit เป็น \tccode{alerts.dismiss}}{unit ผ่าน}{Pass}

  \tccase{2}{follow\_up ปฏิเสธเมื่อยังไม่รับเรื่อง}{%
    \begin{tcsteps}
      \item เรียก \tccode{resolveAlertCloseData} ด้วย \tccode{mode=follow\_up} และ \tccode{alreadyAcknowledged=false}
    \end{tcsteps}}{allowed=false รหัส \tccode{MUST\_ACKNOWLEDGE} ข้อความให้กดรับเรื่องก่อน}{unit ผ่าน}{Pass}

  \tccase{3}{follow\_up ปิดได้หลังรับเรื่องแล้ว}{%
    \begin{tcsteps}
      \item เรียก \tccode{resolveAlertCloseData} ด้วย \tccode{alreadyAcknowledged=true} และเหตุผลปิดเรื่อง
    \end{tcsteps}}{allowed พร้อม \tccode{closedAt} ไม่ทับ \tccode{acknowledgedAt} และ audit เป็น \tccode{alerts.close}}{unit ผ่าน}{Pass}
\end{testscenario}

% TC19 — KKU Gen AI chat provider allowlist
\begin{testscenario}{19}{19}
  \tcmeta{TC19}{CS 2568/01}{ระบบเว็บแอปพลิเคชัน ``NUCARE''}{นายสมชาย ใจดี}{วันที่ 2 พฤศจิกายน พ.ศ.\ 2569}{1.0}
  \tcname{การเชื่อมแชทเพื่อน AI ผ่าน KKU Gen AI โดยจำกัดโมเดล Gemini/GPT/Claude}
  \tcpre{รัน unit test ของ \tccode{kku-gen-models} และ \tccode{kku-gen-provider}}
  \tcdesc{ทดสอบ allowlist โมเดล การเรียก chat completions แบบ OpenAI-compatible และการกันการบ้านโดยไม่ยิง API}

  \tccase{1}{ปฏิเสธโมเดลนอก allowlist}{%
    \begin{tcsteps}
      \item เรียก \tccode{resolveKkuChatModel} ด้วย \tccode{qwen3.7-plus}
    \end{tcsteps}}{ถอยไปโมเดลเริ่มต้นในกลุ่มที่อนุญาต}{unit ผ่าน}{Pass}

  \tccase{2}{บล็อกการบ้านโดยไม่เรียก KKU Gen}{%
    \begin{tcsteps}
      \item สร้าง provider ด้วย fetch จำลอง
      \item ส่งข้อความขอทำการบ้าน
    \end{tcsteps}}{\tccode{blockedReason=homework} และไม่เรียก fetch}{unit ผ่าน}{Pass}

  \tccase{3}{เรียก chat completions เมื่อข้อความปกติ}{%
    \begin{tcsteps}
      \item ส่งข้อความเครียดด้วย persona พี่เมย์
      \item ตรวจ URL Authorization และ system prompt ภาษาไทยเท่านั้น
    \end{tcsteps}}{ได้คำตอบจาก API จำลอง และยิง \tccode{/chat/completions}}{unit ผ่าน}{Pass}
\end{testscenario}

% TC20 — คำขอปรึกษา + นัดหมายผ่าน schedule offer (ที่ปรึกษา / สุขภาพจิต)
\begin{testscenario}{20}{6}
  \tcmeta{TC20}{CS 2568/01}{ระบบเว็บแอปพลิเคชัน ``NUCARE''}{นายสมชาย ใจดี}{วันที่ 5 พฤศจิกายน พ.ศ.\ 2569}{1.3}
  \tcname{คำขอปรึกษานิสิต--ที่ปรึกษา และการนัดปรึกษาสุขภาพจิต}
  \tcpre{เคลียร์ referral ทดสอบแล้ว · มีบัญชีจาก \tccode{seed-local-test-accounts} · API พอร์ต 5000 · รัน \tccode{npx tsx scripts/test-consultation-flow.ts}}
  \tcdesc{ทดสอบหลายเคสจริงผ่าน API: สร้างคำขอ / รับคำขอ / เสนอนัดผ่าน offer แล้วนิสิตยอมรับ / ปิดคำขอ / กันตั้งนัดตรง / กัน IDOR / สร้างนัดสุขภาพจิตหลังรับ referral}

  \tccase{1}{สร้างคำขอปรึกษาหลายเคสและปฏิเสธข้อมูลผิด}{%
    \begin{tcsteps}
      \item ล็อกอินนิสิตทดสอบ 3 คน
      \item POST \tccode{/api/student/consultation-requests} หัวข้อไม่ถูกต้อง · teacherId ที่ไม่มี · ช่องทาง PHONE
      \item สร้างคำขอ 3 รายการไปยัง \tccode{advisor.test} (รับ / นัดผ่าน offer / ปิด)
    \end{tcsteps}}{%
    หัวข้อผิดได้ 400 · teacher ไม่มีได้ 404 · PHONE ได้ 400 · สร้างสำเร็จได้ 201}{%
    ผ่านตามคาด และนิสิตแต่ละคนเห็นคำขอของตน}{Pass}

  \tccase{2}{ที่ปรึกษารับคำขอ เสนอนัดผ่าน offer และปิดคำขอ}{%
    \begin{tcsteps}
      \item GET \tccode{/api/advisor/consultation-requests}
      \item PATCH สถานะ ACCEPTED สำหรับเคส A และ B
      \item PATCH สถานะ SCHEDULED โดยตรง → ต้องถูกปฏิเสธด้วย \tccode{USE\_SCHEDULE\_OFFER}
      \item POST schedule-offers แล้วให้นิสิตยอมรับจนได้ SCHEDULED
      \item PATCH CLOSED สำหรับเคส C และนิสิตยกเลิกเคส A
    \end{tcsteps}}{%
    ตั้งนัดตรงได้ 400 · offer + ยอมรับได้ SCHEDULED · ยกเลิก/ปิดได้ CLOSED}{%
    ผ่าน และนิสิตเห็นสถานะ CLOSED / SCHEDULED / CLOSED ตามลำดับ}{Pass}

  \tccase{3}{กันที่ปรึกษาอื่นจัดการคำขอที่ไม่ใช่ของตน}{%
    \begin{tcsteps}
      \item ล็อกอิน \tccode{advisor.peer}
      \item GET รายการคำขอ และ PATCH รับคำขอของ advisor.test
    \end{tcsteps}}{%
    ไม่เห็นคำขอของคนอื่น และ PATCH ได้ 404}{%
    ผ่าน}{Pass}

  \tccase{4}{รายชื่ออาจารย์คำขอปรึกษาต้องมีที่ปรึกษาที่ผูกนิสิต}{%
    \begin{tcsteps}
      \item GET \tccode{/api/student/consultation-teachers}
      \item ตรวจว่า assigned advisors ถูกดึงก่อน take limit
      \item รันหน่วยทดสอบ \tccode{consultation-teachers.test.ts}
    \end{tcsteps}}{%
    มี \tccode{advisor.test} ในรายการ · หน่วยทดสอบ merge ผ่าน}{%
    ผ่าน}{Pass}

  \tccase{5}{สร้างนัดปรึกษาทีมสุขภาพจิตหลังรับ referral}{%
    \begin{tcsteps}
      \item ที่ปรึกษาส่งต่อ COUNSELING\_CENTER (เหตุผลมี \tccode{[TEST]})
      \item ทีมสุขภาพจิตรับ referral แล้ว POST นัดหมาย
      \item ลองวันเวลาผิดรูป และ studentId ที่ไม่มี
      \item ปิด referral ทดสอบทิ้ง
    \end{tcsteps}}{%
    สร้างนัดได้ 201 · ข้อมูลผิดได้ 400/404 · ปิด referral ได้ 200}{%
    ผ่านครบหลังเคลียร์ข้อมูลเก่า}{Pass}

  \tccase{6}{ผลรวมสคริปต์ integration}{%
    \begin{tcsteps}
      \item รัน \tccode{npx tsx scripts/test-consultation-flow.ts}
    \end{tcsteps}}{%
    PASS ทุกกรณีที่สคริปต์ครอบคลุม}{%
    PASS 64 · FAIL 0}{Pass}
\end{testscenario}

% TC21 — หมุดส่งต่อ + เจรจานัดปรึกษาสองฝ่าย
\begin{testscenario}{21}{5}
  \tcmeta{TC21}{CS 2568/01}{ระบบเว็บแอปพลิเคชัน ``NUCARE''}{นายสมชาย ใจดี}{วันที่ 2 พฤศจิกายน พ.ศ.\ 2569}{1.2}
  \tcname{เส้นทางหมุดติดตามการส่งต่อ และการเจรจานัดปรึกษาจนยินยอมทั้งสองฝ่าย}
  \tcpre{เคลียร์ referral ทดสอบด้วย \tccode{cleanup-test-duplicate-referrals.ts} แล้ว · มีบัญชีทดสอบ · รันหน่วยทดสอบ referral/schedule/active-referral และสคริปต์ flow}
  \tcdesc{ตรวจสอบการสร้างหมุดติดตามเคสส่งต่อ การเจรจานัดปรึกษา และการกันส่งต่อซ้ำเมื่อเคสยังรอรับหรือกำลังดูแล — หลังเคลียร์ข้อมูลปนจากเทสเดิม}

  \tccase{1}{PENDING มีเฉพาะหมุดส่งต่อ}{%
    \begin{tcsteps}
      \item เรียก \tccode{buildReferralRoadmap} ด้วยสถานะ PENDING
      \item ตรวจจำนวนขั้นและ kind ของหมุดแรก
    \end{tcsteps}}{%
    ได้ 1 ขั้น kind=\tccode{sent}}{%
    ผ่านตามหน่วยทดสอบ}{Pass}

  \tccase{2}{ACCEPTED/CLOSED สร้างหมุดครบตามเหตุการณ์}{%
    \begin{tcsteps}
      \item สร้าง roadmap สถานะ ACCEPTED พร้อมบันทึกดูแล
      \item สร้าง roadmap สถานะ CLOSED
    \end{tcsteps}}{%
    ACCEPTED ได้ sent/accepted/care · CLOSED ลงท้ายด้วย closed}{%
    ผ่านตามหน่วยทดสอบ}{Pass}

  \tccase{3}{ข้อเสนอนัดระบุว่าใครต้องตอบ}{%
    \begin{tcsteps}
      \item map ข้อเสนอที่อาจารย์เสนอ
      \item map ข้อเสนอที่นิสิตขอเปลี่ยน
      \item เลือกข้อเสนอ PENDING ล่าสุดด้วย \tccode{pickLatestPendingOffer}
    \end{tcsteps}}{%
    waitingOn เป็น STUDENT / ADVISOR ตามลำดับ และเลือกข้อล่าสุด}{%
    ผ่านตามหน่วยทดสอบ}{Pass}

  \tccase{4}{flow ส่งต่อจริง: สร้าง → รับ → ปิด}{%
    \begin{tcsteps}
      \item รัน \tccode{npx tsx scripts/test-referral-roadmap-flow.ts}
      \item ตรวจ roadmap PENDING/ACCEPTED แล้วให้ MH ปิดเคสทิ้ง
    \end{tcsteps}}{%
    สร้างหมุด sent → accepted สำเร็จ และปิดเคสหลังเทส}{%
    ผ่านครบหลังเคลียร์ DB}{Pass}

  \tccase{5}{กันส่งต่อซ้ำเมื่อเคสกำลังดูแลหรือรอรับ}{%
    \begin{tcsteps}
      \item รัน \tccode{active-referral.test.ts} และ \tccode{test-active-referral-block.ts}
      \item POST ส่งต่อซ้ำเมื่อมีเคส \tccode{PENDING}/\tccode{ACCEPTED}
      \item รัน \tccode{cleanup-test-duplicate-referrals.ts} ลบเหตุผลที่มี \tccode{[TEST]}
    \end{tcsteps}}{%
    หน่วยทดสอบผ่าน · API ได้ 409 · cleanup ลบเคสเทสหมด}{%
    ผ่าน · remaining referrals = none}{Pass}
\end{testscenario}

% TC22 — แยกคิวเจรจานัด / เรียงนัดใกล้ถึง
\begin{testscenario}{22}{4}
  \tcmeta{TC22}{CS 2568/01}{ระบบเว็บแอปพลิเคชัน ``NUCARE''}{นายสมชาย ใจดี}{วันที่ 4 พฤศจิกายน พ.ศ.\ 2569}{1.0}
  \tcname{แยกคิวเจรจานัดของที่ปรึกษา และการเรียงนัดใกล้ถึง}
  \tcpre{มีหน่วยทดสอบ \tccode{consultation-negotiate.test.ts} · รันด้วย \tccode{npx tsx --test src/components/advisor/consultation-negotiate.test.ts}}
  \tcdesc{ตรวจสอบการจัดเลนเจรจานัด (รอฉันตอบ / ยังไม่เสนอ / รอนิสิต) การเรียงลำดับ และการจับนัดภายใน 7 วัน สำหรับ inbox งานวันนี้}

  \tccase{1}{แยกเลนตาม waitingOn ของข้อเสนอ}{%
    \begin{tcsteps}
      \item เรียก \tccode{getNegotiateLane} กับข้อเสนอที่รอ ADVISOR / STUDENT / ไม่มีข้อเสนอ
    \end{tcsteps}}{%
    ได้ waiting\_advisor / waiting\_student / no\_offer ตามลำดับ}{%
    ผ่านตามหน่วยทดสอบ}{Pass}

  \tccase{2}{เรียง ACCEPTED ให้รออาจารย์ขึ้นก่อน}{%
    \begin{tcsteps}
      \item เรียก \tccode{sortAcceptedNegotiations} กับรายการปน 3 เลน
    \end{tcsteps}}{%
    ลำดับเป็นรอฉันตอบ → ยังไม่เสนอ → รอนิสิต}{%
    ผ่าน}{Pass}

  \tccase{3}{เรียง SCHEDULED ตามวันเวลานัด}{%
    \begin{tcsteps}
      \item เรียก \tccode{sortScheduledByAppointment}
    \end{tcsteps}}{%
    นัดที่ใกล้ถึงมาก่อน · ไม่มีวันเวลาอยู่ท้าย}{%
    ผ่าน}{Pass}

  \tccase{4}{นับเลนและนัดใกล้ถึงใน 7 วัน}{%
    \begin{tcsteps}
      \item เรียก \tccode{countNegotiateLanes} และ \tccode{isUpcomingAppointment}
    \end{tcsteps}}{%
    นับถูกต้อง · นัดใน horizon ถูกจับ และนัดไกลเกิน 7 วันไม่ถูกจับ}{%
    ผ่าน · ชุดทดสอบ 7 เคสรวม referral-active ผ่านทั้งหมด}{Pass}
\end{testscenario}

% TC23 — inbox คำขอนัดฝั่งนิสิต
\begin{testscenario}{23}{6}
  \tcmeta{TC23}{CS 2568/01}{ระบบเว็บแอปพลิเคชัน ``NUCARE''}{นายสมชาย ใจดี}{วันที่ 4 พฤศจิกายน พ.ศ.\ 2569}{1.1}
  \tcname{กล่องงานนิสิต: ต้องตอบ ข้อเสนอนัด และนัดใกล้ถึง}
  \tcpre{มีหน่วยทดสอบ \tccode{student-consultation-inbox.test.ts} · รันด้วย \tccode{npx tsx --test src/lib/student-consultation-inbox.test.ts} จากโฟลเดอร์ \tccode{client/}}
  \tcdesc{ตรวจสอบการแยกรายการที่นิสิตต้องตอบ การจับนัดภายใน 7 วัน การนับคิวแบบ actionable และการซ่อนข้อความชุด demo ก่อนแสดงผล}

  \tccase{1}{จับข้อเสนอที่รอให้นิสิตตอบ}{%
    \begin{tcsteps}
      \item เรียก \tccode{isAwaitingStudentResponse} กับข้อเสนอ waitingOn = STUDENT / ADVISOR
    \end{tcsteps}}{%
    ได้ true เฉพาะเมื่อรอ STUDENT}{%
    ผ่านตามหน่วยทดสอบ}{Pass}

  \tccase{2}{เรียงรายการต้องตอบใหม่อยู่ก่อน}{%
    \begin{tcsteps}
      \item เรียก \tccode{listAwaitingStudentResponses} กับรายการปนสถานะ
    \end{tcsteps}}{%
    เหลือเฉพาะที่ต้องตอบ · เรียง createdAt ของข้อเสนอใหม่ก่อน}{%
    ผ่าน}{Pass}

  \tccase{3}{จับและเรียงนัดใกล้ถึงใน 7 วัน}{%
    \begin{tcsteps}
      \item เรียก \tccode{isUpcomingStudentConsultation} และ \tccode{listUpcomingConsultations}
    \end{tcsteps}}{%
    นัดใน horizon ถูกจับ · เรียงใกล้ถึงก่อน · นัดไกลเกิน 7 วันไม่ถูกจับ}{%
    ผ่าน}{Pass}

  \tccase{4}{นับคิว actionable}{%
    \begin{tcsteps}
      \item เรียก \tccode{countStudentActionQueue} กับ PENDING / ต้องตอบ / นัดใกล้ถึง / นัดเก่า / CLOSED
    \end{tcsteps}}{%
    ได้ awaitingStudent=1 · waitingAdvisor=1 · upcoming=1}{%
    ผ่าน}{Pass}

  \tccase{5}{สถานที่นัดและซ่อน DEMO-APPT}{%
    \begin{tcsteps}
      \item เรียก \tccode{resolveAppointmentLocation} และ \tccode{sanitizeStudentFacingNote}
    \end{tcsteps}}{%
    ใช้ scheduleLocation ก่อน · ข้อความที่มีแค่ DEMO-APPT กลายเป็น null}{%
    ผ่าน}{Pass}

  \tccase{6}{ตรวจคำขอเปิดทีละรอบบน helper ฝั่ง client}{%
    \begin{tcsteps}
      \item เรียก \tccode{hasOpenStudentConsultation}
    \end{tcsteps}}{%
    มี SCHEDULED แล้วถือว่ายังเปิด · มีแค่ CLOSED ถือว่าว่าง}{%
    ผ่าน · ชุดทดสอบ 6 เคสผ่านทั้งหมด}{Pass}
\end{testscenario}

% TC24 — คำขอปรึกษาเปิดได้ทีละรอบ
\begin{testscenario}{24}{3}
  \tcmeta{TC24}{CS 2568/01}{ระบบเว็บแอปพลิเคชัน ``NUCARE''}{นายสมชาย ใจดี}{วันที่ 4 พฤศจิกายน พ.ศ.\ 2569}{1.0}
  \tcname{กันนิสิตมีคำขอปรึกษาเปิดพร้อมกันหลายรอบ}
  \tcpre{มีหน่วยทดสอบ \tccode{active-consultation.test.ts} · รันด้วย \tccode{npm run test:unit} จากโฟลเดอร์ \tccode{server/} · API ทดสอบด้วย \tccode{npx tsx scripts/test-consultation-flow.ts}}
  \tcdesc{ตรวจสอบว่าสถานะ PENDING/ACCEPTED/SCHEDULED ถือว่ายังไม่จบ และ API ปฏิเสธการสร้างคำขอใหม่ด้วย 409 จนกว่าจะปิดรอบเดิม}

  \tccase{1}{สถานะที่ถือว่ายังเปิดอยู่}{%
    \begin{tcsteps}
      \item เรียก \tccode{isActiveConsultationStatus} กับ PENDING / ACCEPTED / SCHEDULED / CLOSED
    \end{tcsteps}}{%
    ได้ true เฉพาะสามสถานะแรก}{%
    ผ่านตามหน่วยทดสอบ}{Pass}

  \tccase{2}{ข้อความขัดแย้งแยกตามสถานะ}{%
    \begin{tcsteps}
      \item เรียก \tccode{activeConsultationConflictMessage}
    \end{tcsteps}}{%
    ข้อความระบุว่ามีนัด/คำขอค้างตามสถานะ}{%
    ผ่าน}{Pass}

  \tccase{3}{API ปฏิเสธคำขอซ้ำขณะยังเปิด}{%
    \begin{tcsteps}
      \item ปิดคำขอเปิดของนิสิตทดสอบก่อนสร้าง
      \item สร้างคำขอหนึ่งรายการแล้ว POST ซ้ำทันที
      \item ปิดเคสเทสหลังจบเพื่อไม่สะสมบน dashboard
    \end{tcsteps}}{%
    คำขอที่สองได้ 409 \tccode{ACTIVE\_CONSULTATION} · dashboard ไม่เหลือนัดเทสค้างหลายรายการ}{%
    ผ่าน}{Pass}
\end{testscenario}

% TC25 — ช่องทางออนไลน์/พบตัว + ยกเลิกสองทาง + อาจารย์ระบุสถานที่ตอนเสนอ
\begin{testscenario}{25}{5}
  \tcmeta{TC25}{CS 2568/01}{ระบบเว็บแอปพลิเคชัน ``NUCARE''}{นายสมชาย ใจดี}{วันที่ 8 พฤศจิกายน พ.ศ.\ 2569}{1.4}
  \tcname{รูปแบบคำขอปรึกษาออนไลน์/พบตัว การยกเลิกโดยนิสิตหรืออาจารย์ และการระบุสถานที่โดยอาจารย์}
  \tcpre{มีบัญชีทดสอบ · รันหน่วยทดสอบ \tccode{consultation-channels.test.ts} \tccode{consultation-close.test.ts} และสคริปต์ \tccode{test-consultation-flow.ts}}
  \tcdesc{ตรวจว่านิสิตเลือกได้เฉพาะ ONLINE/IN\_PERSON · อาจารย์เป็นผู้กรอกห้อง/ลิงก์ตอนเสนอนัด · นิสิตขอยกเลิกแล้วอาจารย์ยืนยัน · อาจารย์ยกเลิกเองได้พร้อมเหตุผล (อื่นๆ ต้องกรอกเพิ่ม) · ป้าย ``ยกเลิก'' สีเหลือง}

  \tccase{1}{จำกัดช่องทางที่นิสิตเลือกตอนสร้างคำขอ}{%
    \begin{tcsteps}
      \item รัน \tccode{consultation-channels.test.ts}
      \item POST คำขอด้วย \tccode{preferredChannel=PHONE}
      \item GET ตัวเลือกช่องทางจาก \tccode{/api/student/consultation-teachers}
    \end{tcsteps}}{%
    หน่วยทดสอบผ่าน · PHONE ได้ 400 · ตัวเลือกมีเฉพาะออนไลน์กับพบตัว}{%
    ผ่าน}{Pass}

  \tccase{2}{อาจารย์ระบุสถานที่/ลิงก์ตอนเสนอนัด}{%
    \begin{tcsteps}
      \item รับคำขอแล้ว POST schedule-offers พร้อม location ที่อาจารย์กำหนด
      \item นิสิตยอมรับข้อเสนอ
    \end{tcsteps}}{%
    ได้ SCHEDULED และ \tccode{scheduleLocation} ตามที่อาจารย์เสนอ}{%
    ผ่าน}{Pass}

  \tccase{3}{นิสิตขอยกเลิกแล้วอาจารย์ยืนยัน หรืออาจารย์ยกเลิกเอง}{%
    \begin{tcsteps}
      \item สร้างคำขอแล้วให้อาจารย์รับ (ACCEPTED)
      \item POST \tccode{/api/student/consultation-requests/:id/cancel} พร้อม \tccode{reason}
      \item อาจารย์ PATCH เป็น \tccode{CANCELLED}
      \item ลองยกเลิกเองจาก PENDING ด้วย \tccode{cancelReason} และกรณี \tccode{อื่นๆ} โดยไม่กรอก detail
    \end{tcsteps}}{%
    ได้ \tccode{CANCEL\_REQUESTED} แล้ว \tccode{CANCELLED} · อาจารย์ยกเลิกตรงได้เมื่อมีเหตุผล · อื่นๆ โดยไม่มี detail ได้ 400 · ป้าย ``ยกเลิก'' โทนเหลือง}{%
    ผ่าน}{Pass}

  \tccase{4}{หน่วยทดสอบกฎช่องทางและสิทธิ์ยกเลิก}{%
    \begin{tcsteps}
      \item รัน \tccode{npm run test:unit} ครอบ \tccode{consultation-channels.test.ts}
    \end{tcsteps}}{%
    ทุก assertion ผ่าน รวมสิทธิ์ยกเลิกตรงของอาจารย์}{%
    ผ่าน}{Pass}

  \tccase{5}{ป้ายยกเลิก เหตุผลอาจารย์ และปิดเคสแล้ว}{%
    \begin{tcsteps}
      \item รัน \tccode{consultation-close.test.ts} (CC-01 ถึง CC-04)
    \end{tcsteps}}{%
    เหตุผลนิสิต/อาจารย์ครบ · อื่นๆ บังคับรายละเอียด · \tccode{CANCELLED} เป็น ``ยกเลิก'' · \tccode{CLOSED} ปกติเป็น ``ปิดเคสแล้ว''}{%
    ผ่าน}{Pass}
\end{testscenario}

% TC26 — SOS ฝั่งนิสิตไม่ปิด dialog เมื่อล้มเหลว
\begin{testscenario}{26}{3}
  \tcmeta{TC26}{CS 2568/01}{ระบบเว็บแอปพลิเคชัน ``NUCARE''}{นายสมชาย ใจดี}{วันที่ 5 พฤศจิกายน พ.ศ.\ 2569}{1.0}
  \tcname{การแจ้งช่วยเหลือด่วน (SOS) ของนิสิตเมื่อส่งไม่สำเร็จ}
  \tcpre{มีบัญชีนิสิตทดสอบ · รัน \tccode{student-sos.test.ts}}
  \tcdesc{ตรวจว่าการส่ง SOS ที่ล้มเหลวไม่ถูกมองว่าสำเร็จ และมีทางโทร 1323 / เข้าถึงจากมือถือและแชท}

  \tccase{1}{helper ไม่กลืน error}{%
    \begin{tcsteps}
      \item รัน \tccode{npx tsx --test src/lib/student-sos.test.ts}
    \end{tcsteps}}{%
    สำเร็จคืน ok · ล้มเหลวคืนข้อความ}{%
    PASS 2}{Pass}

  \tccase{2}{dialog คงอยู่เมื่อส่งไม่สำเร็จ}{%
    \begin{tcsteps}
      \item เปิดยืนยัน SOS แล้วจำลอง API error
    \end{tcsteps}}{%
    ไม่ปิด dialog · แสดงข้อความและชี้ 1323}{%
    การ์ดใช้ \tccode{runStudentSosRequest} และปิดเฉพาะเมื่อ ok}{Pass}

  \tccase{3}{ทางเข้าช่วยเหลือบนมือถือและแชทวิกฤต}{%
    \begin{tcsteps}
      \item ตรวจ bottom nav มีลิงก์ \tccode{/student\#emergency-help}
      \item ตรวจห้องแชทเมื่อ crisisSuggested มีปุ่มแจ้งทีมและโทร 1323
    \end{tcsteps}}{%
    เข้าถึงได้โดยไม่ต้องออกจากบริบทวิกฤต}{%
    ปรับใน \tccode{student-mobile-bottom-nav} และ \tccode{chat/page}}{Pass}
\end{testscenario}

\begin{testscenario}{27}{3}
  \tcmeta{TC27}{CS 2568/01}{ระบบเว็บแอปพลิเคชัน ``NUCARE''}{นายสมชาย ใจดี}{วันที่ 5 พฤศจิกายน พ.ศ.\ 2568}{1.0}
  \tcname{การแบ่งหน้ารายการยาวแบบเดียวกับหน้า admin}
  \tcpre{มีชุด helper \tccode{list-pagination} และหน้าที่มีรายการจำนวนมาก เช่น ประวัติคำขอปรึกษาของนิสิต / แท็บปิดแล้วของที่ปรึกษา}
  \tcdesc{ตรวจสอบว่าการตัดหน้าและคำนวณจำนวนแถวต่อหน้าทำงานถูกต้อง ไม่เทรายการทั้งหมดลงหน้าจอ}

  \tccase{1}{ตัดหน้ารายการและแสดงช่วง X--Y จาก Z}{%
    \begin{tcsteps}
      \item เรียก \tccode{slicePageItems} กับรายการ 25 ชิ้น ขนาดหน้า 5
      \item ตรวจหน้า 1 และหน้าสุดท้าย
      \item ตรวจกรณีหน้าเกินจำนวนหน้าและรายการว่าง
    \end{tcsteps}}{คืนช่วงรายการและ \tccode{totalPages} ถูกต้อง พร้อม clamp หน้า}{รัน \tccode{npx tsx --test src/lib/list-pagination.test.ts} ผ่านทั้งชุด \tccode{slicePageItems}}{Pass}

  \tccase{2}{คำนวณแถวต่อหน้าตามขนาดจอ}{%
    \begin{tcsteps}
      \item เรียก \tccode{computeAdaptiveRowsPerPage} สำหรับเดสก์ท็อปและมือถือ
      \item ตรวจค่าอยู่ในช่วง min--max และไม่เกิน max เมื่อจอสูงมาก
    \end{tcsteps}}{จำนวนแถวอยู่ในช่วงที่กำหนดตาม viewport}{เคสในชุด \tccode{computeAdaptiveRowsPerPage} ผ่านครบ}{Pass}

  \tccase{3}{UI รายการยาวใช้ footer ก่อนหน้า/ถัดไป}{%
    \begin{tcsteps}
      \item เปิด \tccode{/student/advisor} ด้วยบัญชีที่มีประวัติจำนวนมาก
      \item ตรวจก้อนประวัติคำขอปรึกษาว่ามีข้อความแสดงช่วงรายการและปุ่มแบ่งหน้า
      \item เปิดคิวคำขอปรึกษาฝั่งที่ปรึกษาแท็บปิดแล้วแล้วสลับหน้า
    \end{tcsteps}}{ไม่เทรายการทั้งหมดในครั้งเดียว และเปลี่ยนหน้าได้}{หน้านิสิตและที่ปรึกษาใช้ \tccode{ListPaginationFooter} เมื่อจำนวนเกินขนาดหน้า}{Pass}
\end{testscenario}

\begin{testscenario}{28}{5}
  \tcmeta{TC28}{CS 2568/01}{ระบบเว็บแอปพลิเคชัน ``NUCARE''}{นายสมชาย ใจดี}{วันที่ 5 พฤศจิกายน พ.ศ.\ 2568}{1.1}
  \tcname{ปรับพอร์ทัลผู้บริหารให้ตรง PDPA และการใช้งาน}
  \tcpre{มีบัญชี EXECUTIVE และ API \tccode{/api/executive/*} พร้อมข้อมูลสรุป}
  \tcdesc{ตรวจสอบการกรองชั้นปี/แผงรายงานใน URL, soft-error, CSV โทนเบา และรายชื่อนิสิต}

  \tccase{1}{สถิติปฏิบัติการกรองตามชั้นปี}{%
    \begin{tcsteps}
      \item เรียก \tccode{buildStudyYearBucket} พร้อมเคสคนละชั้นปี
      \item ตรวจ \tccode{operations.openReferrals}/\tccode{openSos}/\tccode{caseClosureRate}
    \end{tcsteps}}{นับเฉพาะเคสของชั้นปีที่เลือก}{เคส \tccode{EA-05b} ใน \tccode{executive-analytics.test.ts} ผ่าน}{Pass}

  \tccase{2}{อ่าน/สร้าง path ปี หน้า view และ panel}{%
    \begin{tcsteps}
      \item เรียก \tccode{parseExecutiveYearFilter}/\tccode{parseExecutiveReportsPanel}/\tccode{buildExecutivePath}
      \item ตรวจค่า \tccode{?year=} \tccode{?view=} และ \tccode{?panel=comparison}
    \end{tcsteps}}{ปี 1--4 ถูกรับ และ path รวม query รวม panel ได้ถูกต้อง}{ชุด \tccode{executive-year-filter.test.ts} ผ่าน}{Pass}

  \tccase{3}{รายชื่อนิสิตทั้งคณะพร้อมค้นหา/กรอง}{%
    \begin{tcsteps}
      \item เปิด \tccode{/executive?view=students}
      \item โหลดรายชื่อโดยไม่ค้นหา แล้วลองกรองความเสี่ยง/ชั้นปี
      \item เปลี่ยนชั้นปีแล้วตรวจว่าการเลือกนิสิตถูกล้าง
    \end{tcsteps}}{ดูรายละเอียดได้ทุกรายการ พร้อม pagination และล้าง selection เมื่อเปลี่ยนปี}{เมนูรายชื่อนิสิต + \tccode{ExecutiveStudentLookup} ล้าง \tccode{selectedId} เมื่อ \tccode{yearFilter} เปลี่ยน}{Pass}

  \tccase{4}{ส่งออก CSV โทนเบาและรายงานแยกชั้น}{%
    \begin{tcsteps}
      \item ตรวจว่าปุ่ม CSV ไม่อยู่คู่ตัวกรองชั้นปีใน toolbar
      \item เปิดรายงานสรุป/เทียบรอบ แล้วใช้ลิงก์ \tccode{ExecutiveQuietExport}
      \item สลับแท็บสรุปรวม / เทียบรอบ ผ่าน URL \tccode{?panel=}
    \end{tcsteps}}{CSV โทนเบาในหัวข้อ section และเทียบรอบไม่ใช้ปุ่มหลักสีน้ำเงิน}{ใช้ \tccode{ExecutiveQuietExport} ทั้งสรุปและเทียบรอบ}{Pass}

  \tccase{5}{soft-error และแนวโน้มไม่พึ่ง hover}{%
    \begin{tcsteps}
      \item เมื่อรีเฟรชล้มเหลวแต่มีข้อมูลเก่า คงแสดงแดชบอร์ดพร้อมแถบเตือน
      \item เปิดแนวโน้มรายเดือน ตรวจตัวเลข Referral/SOS ใต้แท่งกราฟ
      \item กดปุ่มรีเฟรชใน header และตรวจ badge โซนข้อมูลบนมือถือ
    \end{tcsteps}}{ไม่ทับทั้งหน้าเมื่อ error หลังโหลดสำเร็จ และรายละเอียดแนวโน้มมองเห็นได้โดยไม่ต้อง hover}{อัปเดต \tccode{renderDataState} + caption ใต้กราฟ + badge/รีเฟรชใน header}{Pass}
\end{testscenario}

\begin{testscenario}{29}{4}
  \tcmeta{TC29}{CS 2568/01}{ระบบเว็บแอปพลิเคชัน ``NUCARE''}{นายสมชาย ใจดี}{วันที่ 5 พฤศจิกายน พ.ศ.\ 2568}{1.0}
  \tcname{Dossier นิสิตผู้บริหารและความกว้างเลย์เอาต์}
  \tcpre{มีบัญชี EXECUTIVE และนิสิตที่มีคำขอปรึกษา/ส่งต่อ/ประเมิน}
  \tcdesc{ขยายรายละเอียดนิสิตเป็น dossier ระดับผู้บริหาร และจัดความกว้างเนื้อหาให้เท่า Advisor}

  \tccase{1}{รวมจำนวนปรึกษาต่ออาจารย์}{%
    \begin{tcsteps}
      \item เรียก \tccode{aggregateConsultationsByTeacher} ด้วยหลายอาจารย์
      \item ตรวจ \tccode{totalCount}/\tccode{openCount}/\tccode{closedCount}
    \end{tcsteps}}{จัดกลุ่มตามอาจารย์และเรียงจำนวนมากก่อน}{เคส \tccode{ESD-01} ใน \tccode{executive-student-dossier.test.ts} ผ่าน}{Pass}

  \tccase{2}{สรุปส่งต่อและ SOS}{%
    \begin{tcsteps}
      \item เรียก \tccode{summarizeReferrals} และ \tccode{summarizeEmergencies}
      \item ตรวจนับเปิด/ปิด/เร่งด่วน
    \end{tcsteps}}{ตัวเลขและ recent ถูกต้อง}{เคส \tccode{ESD-02}/\tccode{ESD-03} ผ่าน}{Pass}

  \tccase{3}{API detail คืน dossier}{%
    \begin{tcsteps}
      \item เปิด \tccode{/executive?view=students} แล้วเลือกนิสิต
      \item เรียก \tccode{GET /api/executive/students/:studentUserId}
      \item ตรวจแผงปรึกษาต่ออาจารย์ แนวโน้มคะแนน ส่งต่อ SOS และไทม์ไลน์
    \end{tcsteps}}{แสดง dossier ครบโดยไม่มีข้อความแชท}{controller ใช้ \tccode{getExecutiveStudentDossier} และ UI \tccode{ExecutiveStudentLookup} แสดงครบ}{Pass}

  \tccase{4}{ความกว้างเนื้อหาเท่า Advisor}{%
    \begin{tcsteps}
      \item ตรวจ wrapper หลักของหน้าผู้บริหาร
      \item เทียบกับ Advisor ที่ใช้ \tccode{max-w-7xl}
    \end{tcsteps}}{Executive ใช้ \tccode{max-w-7xl} เช่นเดียวกับ Advisor}{แก้จาก \tccode{max-w-6xl} เป็น \tccode{max-w-7xl} ใน \tccode{executive/page.tsx}}{Pass}
\end{testscenario}

\begin{testscenario}{30}{6}
  \tcmeta{TC30}{CS 2568/01}{ระบบเว็บแอปพลิเคชัน ``NUCARE''}{นายสมชาย ใจดี}{วันที่ 6 พฤศจิกายน พ.ศ.\ 2568}{1.1}
  \tcname{เวิร์กโฟลว์ผู้เชี่ยวชาญสุขภาพจิต (รับคิว นัด บันทึกผล ประวัติ)}
  \tcpre{มีบัญชี \tccode{MENTAL_HEALTH} และ referral เป้าหมาย \tccode{COUNSELING_CENTER}}
  \tcdesc{ตรวจลำดับสถานะคิว นัดหมาย บันทึกผล และประวัตินิสิตแบบ master--detail}

  \tccase{1}{ลำดับสถานะ referral}{%
    \begin{tcsteps}
      \item เรียก \tccode{canTransitionMentalHealthReferralStatus}
      \item ตรวจ PENDING$\to$ACCEPTED, PENDING$\to$CLOSED, ACCEPTED$\to$CLOSED
      \item ตรวจห้าม CLOSED$\to$ACCEPTED
    \end{tcsteps}}{เปลี่ยนสถานะได้ตามลำดับคิวเท่านั้น}{เคส \tccode{MHR-01}--\tccode{MHR-04} ผ่าน}{Pass}

  \tccase{2}{สถานะนัดหมาย}{%
    \begin{tcsteps}
      \item เรียก \tccode{canUpdateMentalHealthAppointmentStatus}
      \item ตรวจ SCHEDULED$\to$COMPLETED/NO\_SHOW/CANCELLED
    \end{tcsteps}}{อัปเดตได้เฉพาะนัดที่ยังเปิด}{เคส \tccode{MHA-S01}/\tccode{MHA-S02} ผ่าน}{Pass}

  \tccase{3}{รับคิวและแจ้งที่ปรึกษา}{%
    \begin{tcsteps}
      \item ผู้เชี่ยวชาญ \tccode{PATCH /api/mental-health/referrals/:id/status} เป็น ACCEPTED
      \item ตรวจว่าที่ปรึกษาได้รับแจ้งเตือนติดตาม
    \end{tcsteps}}{รับเคสสำเร็จและแจ้งผู้ส่งต่อ}{controller เรียก \tccode{notifyUser} ถึง advisor}{Pass}

  \tccase{4}{นัดหมายและบันทึกผล}{%
    \begin{tcsteps}
      \item สร้างนัดจากเคสที่รับแล้ว
      \item อัปเดตสถานะนัดผ่าน \tccode{PATCH .../appointments/:id/status}
      \item บันทึกผลการปรึกษาเฉพาะ referral สถานะ ACCEPTED
    \end{tcsteps}}{ครบวงจรนัด$\to$พบ/ไม่มา/ยกเลิก และบันทึกผลได้}{UI ปฏิทินนัดหมาย + care-records บังคับ ACCEPTED}{Pass}

  \tccase{5}{รับทราบแจ้งเตือนคิว (SOS ฝั่งสตาฟพักไว้)}{%
    \begin{tcsteps}
      \item ตรวจว่าพอร์ทัล MH ไม่แสดงรายการ SOS / ไม่มี route เริ่ม--ปิดเคสฉุกเฉิน
      \item กดรับทราบ alert ประเภทส่งต่อผ่าน \tccode{/api/alerts/:id/acknowledge}
      \item ตรวจว่าปุ่ม SOS ฝั่งนิสิตยังบันทึก \tccode{EmergencyCase} ได้
    \end{tcsteps}}{สตาฟไม่จัดการ SOS ชั่วคราว แต่นิสิตกด SOS ได้}{ถอด emergency handlers จาก MH · นิสิตสร้างเคสโดยไม่แจ้งทีม}{Pass}

  \tccase{6}{ประวัติก่อนรับเคสและการแข่งกดรับ}{%
    \begin{tcsteps}
      \item เรียก \tccode{GET /api/mental-health/referrals/:id/history} ขณะสถานะ PENDING
      \item ตรวจว่าไม่ส่ง \tccode{confidentialNotes} ของรอบก่อน
      \item ผู้เชี่ยวชาญสองคนแข่ง \tccode{PATCH} รับเคสพร้อมกัน
      \item ตรวจว่า UI ไม่แสดงรหัสระบบดิบ และ error ใช้ข้อความภาษาคน
    \end{tcsteps}}{ดูประวัติประเมินได้ก่อนรับโดยไม่เปิดบันทึกลับ คนแรกชนะ คนหลังได้ 409 และรีเฟรช}{history ตาม referral + ตัดบันทึกลับตอน PENDING + \tccode{updateMany} กันแข่ง + \tccode{getApiErrorMessage}}{Pass}
\end{testscenario}

\begin{testscenario}{31}{5}
  \tcmeta{TC31}{CS 2568/01}{ระบบเว็บแอปพลิเคชัน ``NUCARE''}{นายสมชาย ใจดี}{วันที่ 6 พฤศจิกายน พ.ศ.\ 2568}{1.0}
  \tcname{เสริมความพร้อมใช้งานจริงของพอร์ทัลผู้เชี่ยวชาญสุขภาพจิต}
  \tcpre{มีบัญชี \tccode{MENTAL_HEALTH} และ referral เป้าหมาย \tccode{COUNSELING_CENTER}}
  \tcdesc{ตรวจความเป็นเจ้าของเคส การตัดบันทึกลับก่อนรับคิว การล้างนัดเฉพาะของตนเอง และคิวเปิดที่ไม่ถูกตัดทิ้ง}

  \tccase{1}{เจ้าของเคสหลังรับ}{%
    \begin{tcsteps}
      \item เรียก \tccode{canMutateAcceptedReferral}
      \item ตรวจว่าเจ้าของและเคสเก่าที่ยังไม่มี \tccode{acceptedBy} แก้ไขได้
      \item ตรวจว่าผู้เชี่ยวชาญคนอื่นถูกปฏิเสธ
    \end{tcsteps}}{เขียน/ปิดเคสได้เฉพาะผู้รับ}{เคส \tccode{MHO-01}/\tccode{MHO-02} ผ่าน}{Pass}

  \tccase{2}{ตัดบันทึกลับตอน PENDING}{%
    \begin{tcsteps}
      \item เรียก \tccode{redactConfidentialCareNotes}
      \item ตรวจว่า \tccode{confidentialNotes} ถูกตั้งเป็น \tccode{null}
    \end{tcsteps}}{คิวก่อนรับไม่เห็นบันทึกลับ}{เคส \tccode{MHO-03} ผ่าน}{Pass}

  \tccase{3}{รับเคสเก็บ \tccode{acceptedByUserId}}{%
    \begin{tcsteps}
      \item \tccode{PATCH} สถานะเป็น ACCEPTED
      \item ตรวจว่าฟิลด์ \tccode{acceptedByUserId} ถูกตั้งเป็นผู้เรียก
      \item ปิดเคสโดยผู้ใช้อื่นต้องได้ 403
    \end{tcsteps}}{มีเจ้าของถาวรหลังรับ}{controller ตั้งค่าตอน claim และตรวจ ownership}{Pass}

  \tccase{4}{ล้างนัดเฉพาะของตนเอง}{%
    \begin{tcsteps}
      \item เรียก \tccode{POST /api/mental-health/appointments/clear-scheduled}
      \item ตรวจว่า \tccode{updateMany} กรองด้วย \tccode{mentalHealthUserId}
      \item UI ต้องยืนยันก่อนล้าง
    \end{tcsteps}}{ไม่ยกเลิกนัดของเจ้าหน้าที่คนอื่น}{scoped clear + confirm ในหน้าตั้งค่า}{Pass}

  \tccase{5}{คิวเปิดไม่ถูกตัดด้วย take รวม}{%
    \begin{tcsteps}
      \item โหลด dashboard ของผู้เชี่ยวชาญ
      \item ตรวจว่า PENDING/ACCEPTED โหลดแยกจาก CLOSED
    \end{tcsteps}}{คิวรอรับไม่หายเพราะเพดานรายการปิดแล้ว}{แยก open/closed query ใน controller}{Pass}
\end{testscenario}

% TC32 — งานพัฒนานิสิต production hardening
\begin{testscenario}{32}{4}
  \tcmeta{TC32}{CS 2568/01}{ระบบเว็บแอปพลิเคชัน ``NUCARE''}{นายสมชาย ใจดี}{วันที่ 6 พฤศจิกายน พ.ศ.\ 2569}{1.0}
  \tcname{การ hardening พอร์ทัลงานพัฒนานิสิตให้ใช้งานจริง}
  \tcpre{มีนิสิตค้างประเมิน (assignment คลื่น OPEN หรือยังไม่มีผล) และบัญชี \tccode{sd.test@nucare.local}}
  \tcdesc{ตรวจสอบว่า cohort แจ้งเตือนตรงกับแดชบอร์ด กลุ่มเสี่ยงไม่ถูกเรียกว่ากิจกรรม Wellness และฟิลเตอร์ชั้นปีใช้ค่าจริง}

  \tccase{1}{นิยามค้างประเมินตรงกันระหว่างแดชบอร์ดกับ reminder}{%
    \begin{tcsteps}
      \item เรียก unit \tccode{isStudentDevelopmentPendingAssessment}
      \item เปิดแดชบอร์ด SD ดูจำนวนค้างประเมิน
      \item ส่งแจ้งเตือนทั้งหมดโดยไม่ระบุ studentUserIds
    \end{tcsteps}}{%
    ส่งได้เฉพาะกลุ่มค้างประเมินตามนิยามเดียวกับแดชบอร์ด}{%
    unit SDP-01..03 ผ่าน และ API ส่งเฉพาะ pending cohort}{Pass}

  \tccase{2}{ปฏิเสธส่ง reminder ให้นิสิตที่ไม่ค้างประเมิน}{%
    \begin{tcsteps}
      \item POST \tccode{/api/student-development/screening-reminders} พร้อม studentUserIds ของนิสิตที่ประเมินแล้ว
    \end{tcsteps}}{%
    API ตอบ 400 ระบุว่าส่งได้เฉพาะนิสิตค้างประเมิน}{%
    ได้ 400 ตาม validation}{Pass}

  \tccase{3}{ป้ายสรุปกลุ่มเสี่ยงไม่สับสนกับกิจกรรม Wellness}{%
    \begin{tcsteps}
      \item เข้าพอร์ทัล SD
      \item เปิดเมนู ``สรุปกลุ่มเสี่ยง''
    \end{tcsteps}}{%
    แสดงจำนวนตามระดับเสี่ยง ไม่ใช่ระบบลงทะเบียนกิจกรรม}{%
    หน้าแสดงสรุปกลุ่มจากผลคัดกรอง}{Pass}

  \tccase{4}{ฟิลเตอร์ชั้นปีใช้ค่าจริงจากข้อมูล}{%
    \begin{tcsteps}
      \item เปิดมุมมองติดตามการคัดกรอง
      \item เลือกชั้นปีจากตัวเลือกที่สร้างจากรายชื่อจริง
    \end{tcsteps}}{%
    รายชื่อกรองตามชั้นปีที่เลือกได้ถูกต้อง}{%
    ฟิลเตอร์ใช้ yearOptions จาก API ไม่ใช่ค่า hardcode ``ปี 1''}{Pass}
\end{testscenario}

% TC33 — SOS มีผู้รับฝั่งสตาฟ
\begin{testscenario}{33}{4}
  \tcmeta{TC33}{CS 2568/01}{ระบบเว็บแอปพลิเคชัน ``NUCARE''}{นายสมชาย ใจดี}{วันที่ 7 พฤศจิกายน พ.ศ.\ 2569}{1.0}
  \tcname{การแจ้งช่วยเหลือด่วน (SOS) มีผู้รับและดำเนินการได้}
  \tcpre{บัญชีนิสิต มีที่ปรึกษา · บัญชี MENTAL\_HEALTH · รัน unit emergency/SOS}
  \tcdesc{ตรวจว่าเมื่อนิสิตกด SOS ระบบแจ้งทีมสุขภาพจิตและที่ปรึกษา และผู้เชี่ยวชาญรับ/ปิดเคสได้}

  \tccase{1}{ลำดับสถานะ SOS}{%
    \begin{tcsteps}
      \item รัน \tccode{emergency-case-status.test.ts}
    \end{tcsteps}}{%
    OPEN$\to$IN\_PROGRESS/RESOLVED ได้ · ปิดแล้วเปลี่ยนต่อไม่ได้}{%
    ECS-01..04 ผ่าน}{Pass}

  \tccase{2}{รวมผู้รับแจ้งเตือน}{%
    \begin{tcsteps}
      \item รัน \tccode{notify-sos-recipients.test.ts}
    \end{tcsteps}}{%
    รวม MH + ที่ปรึกษา และตัดซ้ำ}{%
    SOS-R01 ผ่าน}{Pass}

  \tccase{3}{นิสิตกด SOS แล้วทีมเห็นคิว}{%
    \begin{tcsteps}
      \item เข้าสู่ระบบนิสิตแล้วกด SOS
      \item เข้า \tccode{mh.test} ดูหน้าหลัก
    \end{tcsteps}}{%
    มีแจ้งเตือน SOS\_OPENED และรายการคำขอช่วยเหลือด่วน}{%
    API สร้าง EmergencyCase + alerts และแดชบอร์ดโหลด emergencyCases}{Pass}

  \tccase{4}{รับดำเนินการแล้วปิดเคส}{%
    \begin{tcsteps}
      \item กดรับดำเนินการแล้วปิดเคสจากหน้าหลัก MH
    \end{tcsteps}}{%
    สถานะ IN\_PROGRESS แล้ว RESOLVED · แจ้งนิสิต · ปิด alerts เมื่อปิดเคส}{%
    PATCH emergency-cases/:id/status ทำงานตามลำดับ}{Pass}
\end{testscenario}

% TC34 — คัดกรอง HIGH แจ้งที่ปรึกษา (ไม่สร้างคิว MH อัตโนมัติ)
\begin{testscenario}{34}{3}
  \tcmeta{TC34}{CS 2568/01}{ระบบเว็บแอปพลิเคชัน ``NUCARE''}{นายสมชาย ใจดี}{วันที่ 7 พฤศจิกายน พ.ศ.\ 2569}{1.2}
  \tcname{ผลคัดกรองความเสี่ยงสูงแจ้งที่ปรึกษาทางแอปและอีเมล}
  \tcpre{นิสิตมีที่ปรึกษาหรือไม่มีก็ได้ · รัน unit HRS-N · EMAIL\_PROVIDER พร้อมหรือ console}
  \tcdesc{เมื่อผลแบบประเมินเป็น HIGH ระบบแจ้งที่ปรึกษา (หรือ SD/ADMIN ถ้ายังไม่มอบหมาย) ไม่สร้าง CaseReferral ไป MH อัตโนมัติ}

  \tccase{1}{วางแผนแจ้งที่ปรึกษา}{%
    \begin{tcsteps}
      \item รัน \tccode{high-risk-screening-notify.test.ts}
    \end{tcsteps}}{%
    มีที่ปรึกษา → notify\_advisors · ไม่มี → notify\_staff\_unassigned}{%
    HRS-N01..02 ผ่าน}{Pass}

  \tccase{2}{หลังส่งแบบ HIGH ไม่โผล่คิว MH อัตโนมัติ}{%
    \begin{tcsteps}
      \item นิสิตส่งแบบประเมินได้ผล HIGH
      \item เข้าพอร์ทัล MH ดูคิวรอรับ
    \end{tcsteps}}{%
    ไม่มี referral จากระบบ · ที่ปรึกษาได้รับ alert/อีเมล}{%
    \tccode{notifyHighRiskScreening} ไม่เรียกสร้าง CaseReferral}{Pass}

  \tccase{3}{ไม่สร้าง EmergencyCase จากคัดกรอง}{%
    \begin{tcsteps}
      \item ตรวจคิวช่วยเหลือด่วน (SOS) หลังผล HIGH
    \end{tcsteps}}{%
    ไม่โผล่เป็น SOS ของนิสิต}{%
    screening ไม่สร้าง EmergencyCase}{Pass}
\end{testscenario}

% TC35 — handoff ข้ามบทบาทหลัง hardening
\begin{testscenario}{35}{4}
  \tcmeta{TC35}{CS 2568/01}{ระบบเว็บแอปพลิเคชัน ``NUCARE''}{นายสมชาย ใจดี}{วันที่ 7 พฤศจิกายน พ.ศ.\ 2569}{1.0}
  \tcname{การเชื่อมโยงบทบาทหลังแก้จุดขาด handoff}
  \tcpre{บัญชีทดสอบครบตาม \tccode{data.md} · unit ที่เกี่ยวข้องผ่าน}
  \tcdesc{ตรวจว่า SOS / คัดกรอง HIGH / สัญญาณแชทวิกฤต / แจ้งเตือนนิสิต / referral คนละ target ทำงานเชื่อมกัน}

  \tccase{1}{SOS มีผู้รับ}{%
    \begin{tcsteps}
      \item นิสิตกด SOS · MH รับ/ปิดที่หน้าหลัก
    \end{tcsteps}}{%
    มีคิวช่วยเหลือด่วนและแจ้งที่ปรึกษา}{%
    TC33 ผ่าน}{Pass}

  \tccase{2}{HIGH สร้างคิว counseling}{%
    \begin{tcsteps}
      \item ผลแบบประเมิน HIGH
    \end{tcsteps}}{%
    MH มีคิวรอรับจากระบบ}{%
    TC34 ผ่าน}{Pass}

  \tccase{3}{นิสิตเห็นการแจ้งเตือน}{%
    \begin{tcsteps}
      \item SD ส่ง reminder · เปิดหน้าหลักนิสิต
    \end{tcsteps}}{%
    มีกล่องการแจ้งเตือนและ toast ไม่กลืนรายการเดิม}{%
    inbox + DAI unit ผ่าน}{Pass}

  \tccase{4}{ส่งต่อคนละ target พร้อมกันได้}{%
    \begin{tcsteps}
      \item ที่ปรึกษาส่งสวัสดิการแล้วยังส่งศูนย์ให้คำปรึกษาได้
    \end{tcsteps}}{%
    ไม่ถูกบล็อกข้าม target}{%
    AR-T01/02 ผ่าน}{Pass}
\end{testscenario}

% TC36 — แอดมินจัดการภาพหน้าแรก / หน้านิสิต
\begin{testscenario}{36}{3}
  \tcmeta{TC36}{CS 2568/01}{ระบบเว็บแอปพลิเคชัน ``NUCARE''}{นายสมชาย ใจดี}{วันที่ 7 พฤศจิกายน พ.ศ.\ 2569}{1.0}
  \tcname{การจัดการภาพประชาสัมพันธ์โดยแอดมิน}
  \tcpre{บัญชีแอดมินตาม \tccode{data.md} · รัน \tccode{npm run test:unit} ที่ \tccode{server/} ให้ผ่านเคส SM-*}
  \tcdesc{ตรวจการตรวจสอบชนิด/ขนาดรูป กัน path traversal ของ URL สาธารณะ และการเพิ่มลบผ่าน API ของแอดมิน}

  \tccase{1}{รับเฉพาะ slot และชนิดรูปที่อนุญาต}{%
    \begin{tcsteps}
      \item เรียก \tccode{isSiteMediaSlot} / \tccode{parseImageDataUrl}
      \item ส่ง data URL PNG 1x1 และชนิด GIF / พาธธรรมดา
    \end{tcsteps}}{%
    PNG ผ่าน · GIF และพาธไม่ใช่ data URL ถูกปฏิเสธ}{%
    SM-01/SM-02 ผ่าน}{Pass}

  \tccase{2}{จำกัดขนาดและ sanitize URL}{%
    \begin{tcsteps}
      \item สร้าง buffer ใหญ่กว่า 2~MiB แล้ว parse
      \item เรียก \tccode{toSiteMediaPublicUrl} ด้วยชื่อที่มี \tccode{..}
    \end{tcsteps}}{%
    ไฟล์ใหญ่ถูกปฏิเสธ · URL ใช้เฉพาะ basename}{%
    SM-03/SM-04 ผ่าน}{Pass}

  \tccase{3}{แอดมินเพิ่ม/ลบแล้วหน้าเว็บดึง API}{%
    \begin{tcsteps}
      \item เข้า \tccode{/admin/site-media} เพิ่มภาพช่อง HOME\_HERO / HOME\_GALLERY / STUDENT\_ACTIVITY
      \item เปิดหน้าแรกและหน้านิสิต · ลบรายการหนึ่งแล้วรีเฟรช
    \end{tcsteps}}{%
    แสดงภาพจาก \tccode{/api/uploads/site-media/...} · ลบแล้วไม่แสดง}{%
    ตรวจด้วยมือหลัง deploy โลคัล}{Pass}
\end{testscenario}

% TC37 — toast แบบ action feedback + helper คัด alerts (ไม่โชว์ toast จาก poll)
\begin{testscenario}{37}{5}
  \tcmeta{TC37}{CS 2568/01}{ระบบเว็บแอปพลิเคชัน ``NUCARE''}{นายสมชาย ใจดี}{วันที่ 8 พฤศจิกายน พ.ศ.\ 2569}{1.3}
  \tcname{Toast หลังกดปุ่ม แยกข้อความตามการกระทำ และ helper คัด alerts}
  \tcpre{รัน \tccode{npx tsx --test src/lib/dashboard-alert-ingest.test.ts src/lib/expandable-toast.test.ts} ที่ \tccode{client/}}
  \tcdesc{Toast กลางใช้เฉพาะ feedback หลังกดปุ่ม/บันทึก — หัวข้อเป็นข้อความการกระทำจริง ไม่ห่อด้วย ``อัปเดตข้อมูลแล้ว'' · ไม่ดึง alerts มาโชว์ popup · Expandable ปิดขยายเป็นค่าเริ่มต้น}

  \tccase{1}{hydrate วิกฤตก่อน (helper)}{%
    \begin{tcsteps}
      \item เรียก \tccode{selectOpenAlertsForFirstHydrate} ด้วยรายการรวม SOS
    \end{tcsteps}}{%
    SOS อยู่ลำดับแรก และไม่เกิน max}{%
    DAI-01 ผ่าน}{Pass}

  \tccase{2}{เกิน max เก็บวิกฤต (helper)}{%
    \begin{tcsteps}
      \item เรียก \tccode{prioritizeAndCapAlerts} ด้วย 5 รายการรวม SOS/HIGH
    \end{tcsteps}}{%
    คืน sos, high และรายการทั่วไปหนึ่ง}{%
    DAI-03 ผ่าน}{Pass}

  \tccase{3}{นับถอยหลังและแถบเวลา}{%
    \begin{tcsteps}
      \item เรียก \tccode{formatToastCountdownSeconds} / \tccode{computeToastElapsedProgress}
    \end{tcsteps}}{%
    ปัดวินาทีถูกต้อง · progress ตามเวลาที่ผ่าน}{%
    ET-01/ET-02 ผ่าน}{Pass}

  \tccase{4}{ไม่เติม body และปิดขยายเป็นค่าเริ่มต้น}{%
    \begin{tcsteps}
      \item เรียก \tccode{resolveExpandableToastContent} โดยไม่มี description และกรณีมีรายละเอียด
    \end{tcsteps}}{%
    ไม่เติมข้อความสำรอง · \tccode{expandByDefault=false} · มี body เมื่อส่ง description ที่ไม่ซ้ำหัวข้อ}{%
    ET-04 ผ่าน}{Pass}

  \tccase{5}{หัวข้อ toast ตามการกระทำ ไม่ใช้คำอัปเดตรวม}{%
    \begin{tcsteps}
      \item เรียก \tccode{resolveActionToastPresentation} ด้วยข้อความรับเรื่อง / error
      \item เรียก \tccode{consultationStatusToastMessage} สำหรับ ACCEPTED / CANCELLED / CLOSED
    \end{tcsteps}}{%
    หัวข้อเป็น ``รับเรื่องแล้ว'' โดยตรง · error ใช้หัวข้อดำเนินการไม่สำเร็จ · สถานะปรึกษาแยกข้อความชัด}{%
    ET-05 ผ่าน}{Pass}
\end{testscenario}

% TC38 — SD จัดการภาพประชาสัมพันธ์หน้านิสิต
\begin{testscenario}{38}{3}
  \tcmeta{TC38}{CS 2568/01}{ระบบเว็บแอปพลิเคชัน ``NUCARE''}{นายสมชาย ใจดี}{วันที่ 7 พฤศจิกายน พ.ศ.\ 2569}{1.0}
  \tcname{สิทธิ์งานพัฒนานิสิตจัดการภาพประชาสัมพันธ์หน้านิสิต}
  \tcpre{รัน \tccode{npx tsx --test src/lib/site-media-access.test.ts} ที่ \tccode{server/}}
  \tcdesc{ตรวจว่า SD แก้ได้เฉพาะช่อง STUDENT\_ACTIVITY และบทบาทอื่น/หน้าแรกถูกปฏิเสธ}

  \tccase{1}{ADMIN ทุกช่อง}{%
    \begin{tcsteps}
      \item เรียก \tccode{getManageableSiteMediaSlots(ADMIN)}
    \end{tcsteps}}{%
    ได้ HOME\_HERO / HOME\_GALLERY / STUDENT\_ACTIVITY}{%
    SMA-01 ผ่าน}{Pass}

  \tccase{2}{SD เฉพาะหน้านิสิต}{%
    \begin{tcsteps}
      \item เรียก \tccode{canManageSiteMediaSlot} สำหรับ STUDENT\_ACTIVITY และ HOME\_HERO
    \end{tcsteps}}{%
    อนุญาตเฉพาะ STUDENT\_ACTIVITY}{%
    SMA-02 ผ่าน}{Pass}

  \tccase{3}{บทบาทอื่นปฏิเสธ}{%
    \begin{tcsteps}
      \item เรียกกับ ADVISOR / STUDENT / MENTAL\_HEALTH
    \end{tcsteps}}{%
    ไม่มีช่องที่จัดการได้}{%
    SMA-03 ผ่าน}{Pass}
\end{testscenario}

% TC39 — กัน production / รหัสผ่านชั่วคราว
\begin{testscenario}{39}{4}
  \tcmeta{TC39}{CS 2568/01}{ระบบเว็บแอปพลิเคชัน ``NUCARE''}{นายสมชาย ใจดี}{วันที่ 7 พฤศจิกายน พ.ศ.\ 2569}{1.0}
  \tcname{การ hardening ความพร้อมใช้งานจริง (รหัสผ่าน อีเมล แชท)}
  \tcpre{รัน unit ที่เกี่ยวข้อง · ไม่ต้องมี SMTP/KKU จริงในโหมดทดสอบ}
  \tcdesc{ตรวจว่ารหัสผ่านชั่วคราวสุ่มได้ ไม่คาดเดา และ production ห้าม console email / mock chat โดยไม่ตั้งใจ}

  \tccase{1}{รหัสผ่านชั่วคราวสุ่ม}{%
    \begin{tcsteps}
      \item รัน \tccode{temporary-password.test.ts}
    \end{tcsteps}}{%
    ได้สตริงยาวพอ ไม่ขึ้นต้น Nu และไม่ซ้ำในชุดตัวอย่าง}{%
    TP-01 ผ่าน}{Pass}

  \tccase{2}{กัน console email ใน production}{%
    \begin{tcsteps}
      \item ตั้ง \tccode{NODE\_ENV=production} และ \tccode{EMAIL\_PROVIDER=console}
      \item เรียก \tccode{assertEmailReadyForEnvironment}
    \end{tcsteps}}{%
    throw เว้นแต่ \tccode{ALLOW\_CONSOLE\_EMAIL=true}}{%
    เคสใน \tccode{email.test.ts} ผ่าน}{Pass}

  \tccase{3}{กัน mock chat ใน production}{%
    \begin{tcsteps}
      \item ตั้ง \tccode{NODE\_ENV=production} และ \tccode{CHAT\_PROVIDER=mock}
      \item เรียก \tccode{assertChatReadyForEnvironment}
    \end{tcsteps}}{%
    throw เว้นแต่ \tccode{ALLOW\_MOCK\_CHAT=true}}{%
    CHAT-ENV-01 ผ่าน}{Pass}

  \tccase{4}{บังคับเปลี่ยนรหัสผ่านหลังสร้างบัญชี}{%
    \begin{tcsteps}
      \item สร้าง/นำเข้าผู้ใช้ใหม่
      \item เข้าสู่ระบบด้วยรหัสชั่วคราว
    \end{tcsteps}}{%
    \tccode{mustChangePassword=true} และพาไป \tccode{/change-password}}{%
    sanitizeUser + MustChangePasswordGate}{Pass}
\end{testscenario}

% TC40 — นิสิตขอนัดผู้เชี่ยวชาญ
\begin{testscenario}{40}{3}
  \tcmeta{TC40}{CS 2568/01}{ระบบเว็บแอปพลิเคชัน ``NUCARE''}{นายสมชาย ใจดี}{วันที่ 7 พฤศจิกายน พ.ศ.\ 2569}{1.0}
  \tcname{นิสิตขอนัดผู้เชี่ยวชาญสุขภาพจิตผ่านช่องทางเฉพาะ}
  \tcpre{นิสิตมีที่ปรึกษา · ไม่มี counseling referral เปิด · รัน unit SCR}
  \tcdesc{ทางเข้าศูนย์ให้คำปรึกษาสำหรับนิสิตคือส่งต่อจากที่ปรึกษา หรือนิสิตขอนัดเอง — ไม่ผูก ConsultationRequest ของที่ปรึกษา}

  \tccase{1}{วางแผนสร้างคำขอ}{%
    \begin{tcsteps}
      \item รัน \tccode{student-counseling-request.test.ts}
    \end{tcsteps}}{%
    สร้างเมื่อมีที่ปรึกษา · ปฏิเสธเมื่อไม่มีหรือมีคิวเปิด}{%
    SCR-01..03 ผ่าน}{Pass}

  \tccase{2}{นิสิตส่งคำขอจากหน้า advisor}{%
    \begin{tcsteps}
      \item เปิด \tccode{/student/advisor\#expert-counseling}
      \item กดส่งคำขอพบผู้เชี่ยวชาญ
    \end{tcsteps}}{%
    สร้าง CaseReferral COUNSELING\_CENTER และแจ้งทีม MH (แอป+อีเมล)}{%
    \tccode{POST /api/student/counseling-requests}}{Pass}

  \tccase{3}{ทีม MH รับคิวแล้วสร้างนัดได้}{%
    \begin{tcsteps}
      \item MH รับ referral
      \item สร้าง MentalHealthAppointment
    \end{tcsteps}}{%
    นัดโผล่ในแดชบอร์ดนิสิต}{%
    ใช้เวิร์กโฟลว์ MH เดิมหลังรับคิว}{Pass}
\end{testscenario}

% TC41 — แชทเพื่อน AI: memory สั้น + เวลา + TTL
\begin{testscenario}{41}{5}
  \tcmeta{TC41}{CS 2568/01}{ระบบเว็บแอปพลิเคชัน ``NUCARE''}{นายสมชาย ใจดี}{วันที่ 7 พฤศจิกายน พ.ศ.\ 2569}{1.1}
  \tcname{แชทเพื่อน AI จำต่อเนื่องแบบประหยัดและกันทำการบ้าน}
  \tcpre{รัน unit guardrails / memory · มีตาราง StudentChatMemory}
  \tcdesc{บอทโทนเพื่อน กันแปะโค้ด/โจทย์ จำเรื่องเดิมด้วยสรุปสั้นพร้อมวันที่ และลบของเก่าเกิน 90 วัน}

  \tccase{1}{กันแปะโค้ดโดยไม่มีคำว่าการบ้าน}{%
    \begin{tcsteps}
      \item รัน \tccode{guardrails.test.ts} เคส CG-05
    \end{tcsteps}}{%
    blocked = homework}{%
    CG-05 ผ่าน}{Pass}

  \tccase{2}{memory สั้นมีวันที่และจำกัดขนาด}{%
    \begin{tcsteps}
      \item รัน \tccode{memory.test.ts}
    \end{tcsteps}}{%
    bullet มีวันที่ไทย · ไม่เกิน 6 บรรทัด · ไม่เก็บโค้ด}{%
    MEM-01..02 ผ่าน}{Pass}

  \tccase{3}{ไม่ฟิกประโยคทักทายจาก memory}{%
    \begin{tcsteps}
      \item เรียก \tccode{buildContinuityGreeting} เมื่อมี memory
      \item ตรวจ \tccode{formatMemoryForPrompt} ส่งแค่บริบทเวลา
    \end{tcsteps}}{%
    เปิดแชทใช้ greeting ของ persona · โมเดลได้บริบทเวลาโดยไม่ถูกบังคับประโยค}{%
    MEM-03, MEM-05 ผ่าน}{Pass}

  \tccase{4}{TTL 90 วัน}{%
    \begin{tcsteps}
      \item ตรวจ \tccode{isStudentChatMemoryStale} และ \tccode{purgeStaleStudentChatMemories}
    \end{tcsteps}}{%
    ของเก่าเกิน 90 วันถูกลบตอนโหลด/ตอนสตาร์ทเซิร์ฟเวอร์}{%
    MEM-06 ผ่าน}{Pass}

  \tccase{5}{ไม่เก็บประวัติแชทเต็มในฐานข้อมูล}{%
    \begin{tcsteps}
      \item ตรวจโมเดล \tccode{StudentChatMemory}
    \end{tcsteps}}{%
    มีแค่ summary / lastTopic / turnCount / updatedAt ต่อนิสิต+persona}{%
    ไม่มีตารางข้อความแชทเต็ม}{Pass}
\end{testscenario}

\begin{testscenario}{42}{6}
  \tcmeta{TC42}{CS 2568/01}{ระบบเว็บแอปพลิเคชัน ``NUCARE''}{นายสมชาย ใจดี}{วันที่ 8 พฤศจิกายน พ.ศ.\ 2568}{1.2}
  \tcname{แยกคิวรอรับกับเคสที่ดูแล และปรับ UX triage}
  \tcpre{มีบัญชี \tccode{MENTAL_HEALTH} และ referral สถานะ PENDING/ACCEPTED}
  \tcdesc{ตรวจการแยกเมนูคิว/เคส ข้อความอายุคิว การนำทางจากแดชบอร์ด helper อายุคิว และป้ายสถานะกิจกรรม}

  \tccase{1}{ข้อความอายุคิว}{%
    \begin{tcsteps}
      \item เรียก \tccode{formatQueueWait} ด้วยเวลาอ้างอิงคงที่
      \item ตรวจเคสเพิ่งเข้าคิว / นาที / ชั่วโมง / วัน และวันที่ไม่ถูกต้อง
    \end{tcsteps}}{ได้ข้อความไทยสั้นตามช่วงเวลา}{เคส unit ใน \tccode{client/src/lib/format-queue-wait.test.ts} ผ่าน}{Pass}

  \tccase{2}{เมนูคิวกับเคสแยกกัน}{%
    \begin{tcsteps}
      \item เปิด \tccode{/mental-health?view=queue} เห็นเฉพาะ PENDING
      \item เปิด \tccode{/mental-health?view=cases} เห็นเฉพาะ ACCEPTED
      \item ตรวจรายการคิวแสดงเหตุผลส่งต่อและอายุคิว
    \end{tcsteps}}{คิวกับเคสไม่ปนกัน และคิว triage ได้จากรายการซ้าย}{UI แยก mode \tccode{queue}/\tccode{cases} และแสดง reason + wait}{Pass}

  \tccase{3}{แดชบอร์ดนำทางได้}{%
    \begin{tcsteps}
      \item กด MetricCard ``คิวรอรับ'' / ``เคสที่ดูแล'' / ``นัดหมายวันนี้''
      \item กดแถวพรีวิวคิวหรือหัวข้อเคสบนแดชบอร์ด
    \end{tcsteps}}{พาไปวิวที่ถูกต้อง และ deep-link \tccode{?referral=} เมื่อมี}{MetricCard มี \tccode{onClick}; พรีวิวเรียก \tccode{navigateTo}}{Pass}

  \tccase{4}{โทนเร็วและความลับ}{%
    \begin{tcsteps}
      \item ตรวจว่าบันทึกลับต้องกด ``แสดงบันทึกลับ'' ก่อนเห็นเนื้อหา
      \item ตรวจความเสี่ยง/เร่งด่วนเป็นป้ายข้างชื่อ ไม่ใช่แบนเนอร์สีชมพู
      \item ตรวจแท็บเริ่มต้นของคิวเป็นภาพรวม
    \end{tcsteps}}{ไม่เพิ่มขั้นตอนหนัก แต่ปกป้องบันทึกลับและเน้น triage}{reveal บันทึกลับ + risk/urgent chips + default overview}{Pass}

  \tccase{5}{ป้ายสถานะกิจกรรมล่าสุด}{%
    \begin{tcsteps}
      \item เรียก \tccode{presentActivityStatus} สำหรับ CONSULTATION/REFERRAL/ASSESSMENT
      \item ตรวจ UI กิจกรรมล่าสุดแสดงประเภท สถานะจบ/เปิด ผู้เกี่ยวข้อง
    \end{tcsteps}}{แยกจบกับยังเปิดได้ชัด ไม่เหลือแค่หัวข้อกับวันที่}{unit ผ่าน; รายการโชว์ type + status badge + actor}{Pass}

  \tccase{6}{เลิกใช้ localsay ตอนปิดเคส/ปฏิเสธคิว}{%
    \begin{tcsteps}
      \item กดปิดเคสหรือปฏิเสธคิวในพอร์ทัล MH
      \item ตรวจว่าขึ้น \tccode{ConfirmDialog} ในแอป ไม่ใช่ \tccode{window.confirm}
      \item ค้นทั้ง client ว่าไม่มี \tccode{window.confirm}/\tccode{window.prompt} คงเหลือ
    \end{tcsteps}}{ยืนยันด้วย popup ในระบบเท่านั้น}{MH ใช้ ConfirmDialog; client ไม่เหลือ localsay}{Pass}
\end{testscenario}

\begin{testscenario}{43}{3}
  \tcmeta{TC43}{CS 2568/01}{ระบบเว็บแอปพลิเคชัน ``NUCARE''}{นายสมชาย ใจดี}{วันที่ 9 พฤศจิกายน พ.ศ.\ 2568}{1.1}
  \tcname{เปิดโมดัลแล้วแถบบนไม่หาย และล็อกสกอลล์พื้นหลังอย่างปลอดภัย}
  \tcpre{มีบัญชีอาจารย์ที่ปรึกษา และมีสัญญาณเข้าที่ปิดเคสได้}
  \tcdesc{กันบั๊ก Chrome ที่ backdrop-blur ของโมดัลทับ sticky header ที่มี blur แล้วทำให้แถบบนหาย พร้อมรวมการล็อกสกอลล์เป็น utility เดียว และกันสกอลล์บาร์ซ้อนดันเนื้อหาซ้าย}

  \tccase{1}{ล็อกสกอลล์แบบนับจำนวน}{%
    \begin{tcsteps}
      \item รัน unit \tccode{client/src/lib/body-scroll-lock.test.ts}
      \item ตรวจล็อกครั้งแรกตั้ง \tccode{position: fixed} และคง \tccode{scrollY}
      \item ตรวจล็อกซ้อนกันและปลดซ้ำไม่ทำให้สถานะพัง
      \item ตรวจว่าไม่ใส่ \tccode{paddingRight} ทับ \tccode{scrollbar-gutter: stable}
    \end{tcsteps}}{เทส BSL-01 ถึง BSL-04 ผ่านทั้งหมด}{รัน \tccode{npx tsx --test src/lib/body-scroll-lock.test.ts} ผ่าน 4/4}{Pass}

  \tccase{2}{แถบบนยังเห็นตอนเปิดป๊อปอัป}{%
    \begin{tcsteps}
      \item เปิด \tccode{/advisor?view=academic}
      \item กด ``บันทึกและปิดเคส'' ให้ขึ้น \tccode{PromptDialog}
      \item สังเกตแถบ sticky ด้านบน (ชื่อ/เมนู) ยังอยู่ใต้ overlay
    \end{tcsteps}}{แถบบนไม่หาย ไม่เลื่อนหลุด และหน้าหลังไม่กระตุก}{เอา backdrop-blur ออกจาก overlay; header staff ใช้พื้นขาวทึบ; PromptDialog ใช้ \tccode{useBodyScrollLock}}{Pass}

  \tccase{3}{จุดโมดัลอื่นใช้ล็อกเดียวกัน และไม่ดันเลย์เอาต์}{%
    \begin{tcsteps}
      \item ตรวจ \tccode{ConfirmDialog}/\tccode{StaffModalOverlay}/\tccode{DeleteSurveyDialog}/\tccode{AuthPopup} เรียก utility กลาง
      \item ค้น client ว่าไม่เหลือการตั้ง \tccode{document.body.style.overflow = "hidden"} นอก utility
      \item เปิดรายละเอียดข้อสังเกตแล้วดูว่าข้อความใต้หัวข้อไม่ขยับซ้าย
    \end{tcsteps}}{ล็อกสกอลล์รวมศูนย์ ไม่มีสกอลล์บาร์ซ้อนดันพื้นหลัง}{overlay ใช้ \tccode{overflow-hidden}; lock ไม่ใส่ paddingRight; เหลือ overflow lock เฉพาะ utility}{Pass}
\end{testscenario}

\begin{testscenario}{44}{3}
  \tcmeta{TC44}{CS 2568/01}{ระบบเว็บแอปพลิเคชัน ``NUCARE''}{นายสมชาย ใจดี}{วันที่ 10 พฤศจิกายน พ.ศ.\ 2568}{1.2}
  \tcname{ปฏิทินนัดใกล้ถึง 7 วันบนหน้าแรกที่ปรึกษา}
  \tcpre{มีบัญชีอาจารย์ที่ปรึกษา และมีนัด SCHEDULED ในช่วง 7 วัน (อย่างน้อย 1 รายการ)}
  \tcdesc{ปฏิทินแนว week timeline คอลัมน์วัน × แกนชั่วโมง วางการ์ดนัดตามเวลา}

  \tccase{1}{จัดกลุ่มนัดและโมเดล timeline}{%
    \begin{tcsteps}
      \item รัน unit \tccode{client/src/components/advisor/upcoming-week-calendar.test.ts}
      \item ตรวจคีย์วันที่ จัดกลุ่มรายวัน และตำแหน่งชั่วโมงของนัด
    \end{tcsteps}}{เทส UWC-01 ถึง UWC-04 ผ่าน}{รัน \tccode{npx tsx --test src/components/advisor/upcoming-week-calendar.test.ts} ผ่าน 4/4}{Pass}

  \tccase{2}{UI week view บนงานวันนี้}{%
    \begin{tcsteps}
      \item เปิด \tccode{/advisor} หน้างานวันนี้
      \item เห็นหัวสัปดาห์ คอลัมน์ 7 วัน แกนชั่วโมง และการ์ดนัดบนตาราง
    \end{tcsteps}}{เห็นหลายวันพร้อมกันบน timeline ไม่ต้องกดทีละวัน}{\tccode{UpcomingWeekTimeline} แสดง week grid + event cards}{Pass}

  \tccase{3}{ว่างและลิงก์ดูทั้งหมด}{%
    \begin{tcsteps}
      \item เมื่อไม่มีนัดใน 7 วัน แสดง empty
      \item กด ``ดูทั้งหมด'' ไปแท็บนัดแล้ว
    \end{tcsteps}}{empty เมื่อไม่มีนัด และนำทางไป consultations/scheduled ได้}{empty state; ปุ่มดูทั้งหมดเรียก \tccode{onOpenAll}}{Pass}
\end{testscenario}

\begin{testscenario}{45}{4}
  \tcmeta{TC45}{CS 2568/01}{ระบบเว็บแอปพลิเคชัน ``NUCARE''}{นายสมชาย ใจดี}{วันที่ 10 พฤศจิกายน พ.ศ.\ 2568}{1.0}
  \tcname{จำกัดรายการแจ้งเตือน ปฏิทินหลายวัน และกันชนนัด 1 ชั่วโมง}
  \tcpre{มีบัญชีนิสิตที่มีการแจ้งเตือนหลายรายการ และบัญชีที่ปรึกษาที่มีนัด SCHEDULED อย่างน้อย 1 รายการ}
  \tcdesc{กันรายการแจ้งเตือนยืดไม่หยุด ปรับปฏิทินให้อ่านหลายวันพร้อมกัน และป้องกันเสนอนัดทับช่วง}

  \tccase{1}{รายการแจ้งเตือนมี max-height}{%
    \begin{tcsteps}
      \item เปิดหน้าแรกนิสิตเมื่อมีการแจ้งเตือนหลายรายการ
      \item ตรวจกล่อง ``การแจ้งเตือน'' เลื่อนในกรอบได้ ไม่ยืดยาวทั้งหน้า
    \end{tcsteps}}{รายการอยู่ในกรอบสูงจำกัดและ scroll ได้}{ใช้ max-height + overflow-y-auto ในกล่องการแจ้งเตือน}{Pass}

  \tccase{2}{ปฏิทินเห็นหลายวันพร้อมกัน}{%
    \begin{tcsteps}
      \item เปิด \tccode{/advisor} งานวันนี้
      \item ตรวจแถบ 7 วันและรายการนัดแยกตามวันโดยไม่ต้องกดวันทีละวัน
    \end{tcsteps}}{เห็นนัดทุกวันที่มีในสัปดาห์พร้อมกัน}{agenda จัดกลุ่มตามวันใต้แถบสัปดาห์}{Pass}

  \tccase{3}{กันชนนัด 1 ชั่วโมง (unit)}{%
    \begin{tcsteps}
      \item รัน \tccode{server/.../consultation-schedule-conflict.test.ts}
      \item รัน \tccode{client/.../consultation-schedule-conflict.test.ts}
    \end{tcsteps}}{เทส CSC ผ่านทั้งฝั่ง server และ client}{server 3/3 และ client 2/2 ผ่าน}{Pass}

  \tccase{4}{popup เสนอนัดโชว์นัดที่มีและกันทับ}{%
    \begin{tcsteps}
      \item เปิดเสนอวัน/สถานที่นัดเมื่อมีนัดอื่นอยู่แล้ว
      \item เห็นรายการนัด/ข้อเสนอรอตอบใน popup
      \item ลองเลือกเวลาในชั่วโมงเดียวกับนัดเดิม — ระบบปฏิเสธ
    \end{tcsteps}}{แสดง busy slots และบล็อกทั้ง client + API 409}{ตรวจก่อน submit บน client; server ตรวจตอน create/accept offer}{Pass}
\end{testscenario}

% TC46 — AUTH_COOKIE_SECURE สำหรับ faculty HTTP staging
\begin{testscenario}{46}{4}
  \tcmeta{TC46}{CS 2568/01}{ระบบเว็บแอปพลิเคชัน ``NUCARE''}{นายสมชาย ใจดี}{วันที่ 10 พฤศจิกายน พ.ศ.\ 2569}{1.0}
  \tcname{การตั้งค่า Secure flag ของ cookie ยืนยันตัวตน}
  \tcpre{รัน \tccode{src/lib/auth-cookie.test.ts} · ไม่ต้องมี Docker}
  \tcdesc{ตรวจว่า production ค่าเริ่มต้นใช้ Secure และ faculty HTTP staging ปิด Secure ผ่าน env ได้}

  \tccase{1}{production ไม่ตั้ง env}{%
    \begin{tcsteps}
      \item เรียก \tccode{resolveAuthCookieSecure} ด้วย \tccode{NODE\_ENV=production}
    \end{tcsteps}}{%
    คืนค่า \tccode{true}}{%
    ผ่านตาม unit test}{Pass}

  \tccase{2}{development ไม่ตั้ง env}{%
    \begin{tcsteps}
      \item เรียกด้วย \tccode{NODE\_ENV=development}
    \end{tcsteps}}{%
    คืนค่า \tccode{false}}{%
    ผ่านตาม unit test}{Pass}

  \tccase{3}{faculty HTTP staging}{%
    \begin{tcsteps}
      \item ตั้ง \tccode{NODE\_ENV=production} และ \tccode{AUTH\_COOKIE\_SECURE=false}
    \end{tcsteps}}{%
    คืนค่า \tccode{false} เพื่อให้ login บน \tccode{http://ip:89} ได้}{%
    ผ่านตาม unit test}{Pass}

  \tccase{4}{บังคับ Secure นอก production}{%
    \begin{tcsteps}
      \item ตั้ง \tccode{AUTH\_COOKIE\_SECURE=true} ใน development
    \end{tcsteps}}{%
    คืนค่า \tccode{true}}{%
    ผ่านตาม unit test}{Pass}
\end{testscenario}


% แม่แบบว่าง: คัดลอก scenarios/_template.tex เป็นไฟล์ใหม่ แล้ว \input ที่นี่

\end{document}
